% ── chemistry-handout-preamble.tex ── 化学竞赛讲义 LaTeX 导言区（v7 · 精排版）
\documentclass[a4paper,11pt]{ctexart}

% ── 字体优化 ──
\setCJKmainfont[AutoFakeBold=true,BoldFont=SimHei]{SimSun}
\setCJKsansfont{SimHei}

% 化学式宏（mhchem）
\usepackage[version=4]{mhchem}

% ── 数学 ──
\usepackage{amsmath,amssymb}

% ── 图片 ──
\usepackage{graphicx}
\usepackage[font=small,labelfont=bf,skip=4pt]{caption}

% ── 表格（跨页+自动列宽）──
\usepackage{booktabs}
\usepackage{xltabular}
\usepackage{array}
\newcolumntype{L}{>{\raggedright\arraybackslash}X}
\newcolumntype{C}{>{\centering\arraybackslash}X}
\newcolumntype{R}{>{\raggedleft\arraybackslash}X}

% ── 浮动体（[H] 强制定位）──
\usepackage{float}

% ── 页面 ──
\usepackage[includeheadfoot,margin=2.2cm,headheight=14pt]{geometry}
\setlength{\textwidth}{16.6cm}

% ── 行间距 ──
\usepackage{setspace}
\setstretch{1.15}
\setlength{\parskip}{0.3em}
\setlength{\parindent}{2em}

% ── 防止表格跨页 ──
\usepackage{longtable}
\setlength{\LTpre}{0pt}
\setlength{\LTpost}{0pt}

% ── 防止图片跨页（尽量）──
\renewcommand{\textfraction}{0.1}
\renewcommand{\topfraction}{0.9}
\renewcommand{\bottomfraction}{0.9}
\setcounter{totalnumber}{10}
\setcounter{topnumber}{10}
\setcounter{bottomnumber}{10}
\setlength{\textfloatsep}{10pt plus 2pt minus 2pt}
\setlength{\floatsep}{8pt plus 2pt minus 2pt}
\setlength{\intextsep}{10pt plus 2pt minus 2pt}

% ── 超链接(无红框) ──
\usepackage[colorlinks=true,linkcolor=black,citecolor=black,urlcolor=black]{hyperref}

% ── 颜色 ──
\usepackage[dvipsnames,svgnames]{xcolor}
\definecolor{chapterblue}{RGB}{23,50,77}
\definecolor{insightbg}{RGB}{235,240,250}
\definecolor{warningbg}{RGB}{255,245,235}
\definecolor{tipbg}{RGB}{235,250,235}

% ── 页眉页脚 ──
\usepackage{fancyhdr}
\pagestyle{fancy}
\fancyhf{}
\fancyhead[L]{\small\color{chapterblue}\nouppercase{\leftmark}}
\fancyhead[R]{\small\color{chapterblue}\thepage}
\renewcommand{\headrule}{{\color{chapterblue}\hrule height 0.8pt}}
% 去掉页眉中的"超级充实版"等后缀
\renewcommand{\sectionmark}[1]{\markboth{#1}{}}
\patchcmd{\section}{\markright{\thesection\ --\ #1}}{\markright{\thesection}}{}{}

% ── 章节编号与样式 ──
\setcounter{secnumdepth}{1}
\ctexset{
  section = {
    format = \Large\bfseries\heiti\color{chapterblue},
    aftername = \hspace{0.5em},
    beforeskip = 2ex,
    afterskip = 1ex,
  },
  subsection = {
    format = \large\bfseries\heiti,
    beforeskip = 1.5ex,
    afterskip = 0.5ex,
  },
  subsubsection = {
    format = \normalsize\bfseries\heiti,
    beforeskip = 1ex,
    afterskip = 0.3ex,
  },
  paragraph = {
    format = \small\bfseries,
    beforeskip = 0.5ex,
    afterskip = 0.2ex,
  },
}

% ── blockquote（引用/提示框）缩减 ──
\usepackage{etoolbox}
\AtBeginEnvironment{quote}{\small\setstretch{1.05}\leftskip=1em\rightskip=1em}

% ── 教学提示框（浅底色版）──
\newenvironment{insightbox}{%
  \par\smallskip
  \noindent{\small\bfseries\color{RoyalBlue}🧠 认知冲突}\par
  \leavevmode\colorbox{insightbg}{\parbox{\dimexpr\textwidth-2\fboxsep\relax}{%
  \ignorespaces}}%
}{\par\smallskip}

\newenvironment{warningbox}{%
  \par\smallskip
  \noindent{\small\bfseries\color{BrickRed}⚠️ 易错点}\par
  \leavevmode\colorbox{warningbg}{\parbox{\dimexpr\textwidth-2\fboxsep\relax}{%
  \ignorespaces}}%
}{\par\smallskip}

\newenvironment{tipbox}{%
  \par\smallskip
  \noindent{\small\bfseries\color{ForestGreen}💡 理解提示}\par
  \leavevmode\colorbox{tipbg}{\parbox{\dimexpr\textwidth-2\fboxsep\relax}{%
  \ignorespaces}}%
}{\par\smallskip}

% ── 列表（答案区加间距）──
\usepackage{enumitem}
\setlist{nosep}
% 专门用于参考答案的列表（每项之间空一行）
\newlist{answerlist}{enumerate}{1}
\setlist[answerlist]{label=\arabic*., leftmargin=2em, itemsep=0.8em}

% ── 防止内容溢出 ──
\usepackage{microtype}
\setlength{\emergencystretch}{2em}
\setlength{\hfuzz}{2pt}
\setlength{\overfullrule}{0pt}
\sloppy

% ── pandoc 兼容 ──
\providecommand{\tightlist}{\setlength{\itemsep}{0pt}\setlength{\parskip}{0pt}}
\newenvironment{Shaded}{}{}
\newenvironment{Highlighting}{}{}
\newcommand{\NormalTok}[1]{#1}
\newcommand{\CommentTok}[1]{#1}
\newcommand{\KeywordTok}[1]{#1}
\newcommand{\StringTok}[1]{#1}
\newcommand{\FunctionTok}[1]{#1}
\newcommand{\DataTypeTok}[1]{#1}
\newcommand{\NumberTok}[1]{#1}
\newcommand{\ControlFlowTok}[1]{#1}
\newcommand{\OperatorTok}[1]{#1}
\newcommand{\BaseNTok}[1]{#1}
\newcounter{none}

% ── 图片尺寸限制 ──
\AtBeginDocument{\setkeys{Gin}{width=\textwidth,keepaspectratio}}

% ── 防止图片跨节 ──
\usepackage{placeins}

% ── 自定义命令 ──
% 练习题标题（去掉括号）
% 例题标题（去掉中括号）
% 这些需要在convert.py预处理中做正则替换

% ── Unicode 符号渲染命令（preprint 中调用，\ensuremath 避免 $...$ 切断 CJK 字间距）──
\newcommand{\箭}{\ensuremath{\rightarrow}}
\newcommand{\左箭}{\ensuremath{\leftarrow}}
\newcommand{\上箭}{\ensuremath{\uparrow}}
\newcommand{\下箭}{\ensuremath{\downarrow}}
\newcommand{\衡}{\ensuremath{\rightleftharpoons}}
\newcommand{\双箭}{\ensuremath{\leftrightarrow}}
\newcommand{\约等于}{\ensuremath{\approx}}
\newcommand{\小于等于}{\ensuremath{\leq}}
\newcommand{\大于等于}{\ensuremath{\geq}}
\newcommand{\不等于}{\ensuremath{\neq}}
\newcommand{\西格马}{\ensuremath{\sigma}}
\newcommand{\派}{\ensuremath{\pi}}

\graphicspath{{C:/Obsidion/妙妙屋/11-模板/scripts/.chem_media/}}
\title{晶体学与晶体结构}
\date{}
\setlength{\tabcolsep}{3.5pt}
\small
\begin{document}
\thispagestyle{plain}
\begin{center}
{\zihao{2}\bfseries 晶体学与晶体结构\par}
\end{center}
\vspace{1em}
\section{晶体学与晶体结构}\label{ux6676ux4f53ux5b66ux4e0eux6676ux4f53ux7ed3ux6784}

\begin{center}\rule{0.5\linewidth}{0.5pt}\end{center}

\subsection{学习目标}\label{ux5b66ux4e60ux76eeux6807}

\begin{itemize}
\tightlist
\item
  了解晶体的四个基本性质(自范性、各向异性、对称性、均一性)
\item
  能理解点阵与结构基元的关系(晶体结构 = 点阵 + 结构基元)，能辨认简单晶体的结构基元
\item
  了解七大晶系和14种布拉维格子的基本特征
\item
  能计算晶胞中的原子数(顶角 \(1/8\)、棱心 \(1/4\)、面心 \(1/2\)、体内 \(1\))
\item
  了解素晶胞与复晶胞(P、I、F、C)，并能用平移操作检验晶胞类型
\item
  能用分数坐标描述晶胞中原子的位置
\item
  能识别分子/晶体中的基本对称元素(旋转轴 \(C_n\)、镜面 \(\sigma\)、对称中心 \(i\))
\end{itemize}

\begin{center}\rule{0.5\linewidth}{0.5pt}\end{center}

\subsection{一、晶体与晶体的基本性质}\label{ux4e00ux6676ux4f53ux4e0eux6676ux4f53ux7684ux57faux672cux6027ux8d28}

\subsubsection{1.1 晶体的定义}\label{ux6676ux4f53ux7684ux5b9aux4e49}

晶体是原子、离子或分子按照一定的周期性在空间排列，形成具有一定规则几何外形的固体。与非晶态物质(玻璃、树脂)相比，晶体具有长程有序的微观结构。

\begin{quote}
\textbf{深度辨析}：「液晶 (Liquid Crystal)」---介于晶体与液体之间的中间态。液晶分子具有取向有序(像晶体)但位置无序(像液体)。液晶的发现挑战了''晶体 vs 非晶体''的二元分类，也因此获得了1991年诺贝尔物理学奖(de Gennes)。液晶显示器(LCD)正是利用液晶的电光效应工作的。
\end{quote}

\subsubsection{1.2 晶体的四个基本性质}\label{ux6676ux4f53ux7684ux56dbux4e2aux57faux672cux6027ux8d28}

{\def\LTcaptype{none} % do not increment counter
\begin{longtable}[]{@{}
  >{\raggedright\arraybackslash}p{(\linewidth - 4\tabcolsep) * \real{0.3333}}
  >{\raggedright\arraybackslash}p{(\linewidth - 4\tabcolsep) * \real{0.3333}}
  >{\raggedright\arraybackslash}p{(\linewidth - 4\tabcolsep) * \real{0.3333}}@{}}
\toprule\noalign{}
\begin{minipage}[b]{\linewidth}\raggedright
性质
\end{minipage} & \begin{minipage}[b]{\linewidth}\raggedright
定义
\end{minipage} & \begin{minipage}[b]{\linewidth}\raggedright
实例说明
\end{minipage} \\
\midrule\noalign{}
\endhead
\bottomrule\noalign{}
\endlastfoot
\textbf{自范性} & 在适宜条件下，晶体能自发地呈现封闭的、规则的凸多面体外形 & 氯化钠晶体呈立方体，石英晶体呈六方柱 \\
\textbf{各向异性} & 沿晶体的不同方向，物理性质(折射率、热导率、解理性等)不同 & 云母沿层状方向易剥离，垂直方向难；石墨沿层方向导电，垂直方向不导电 \\
\end{longtable}
}

\begin{quote}
\textbf{认知冲突}：「玻璃是固体，为什么不是晶体？」---因为玻璃中的原子/分子排列是\textbf{短程有序但长程无序}(像冻结的液体)。它没有固定的熔点，只是逐渐软化。这揭示了晶体的本质特征不是''硬''或''固体''，而是\textbf{周期性长程有序}。
\end{quote}

\textbf{对称性} \textbar{} 晶体的理想外形和内部结构具有对称性 \textbar{} 立方晶体绕 \(C_4\) 轴旋转 \(90^\circ\) 后与自身重合 \textbar{}\\
\textbf{均一性} \textbar{} 同一晶体的各个部分宏观性质相同 \textbar{} 晶体内任意取两点，化学组成和物理性质一致 \textbar{}

\begin{quote}
\textbf{注意}：均一性并不与各向异性矛盾---均一性指同一方向上的性质各处相同，各向异性指不同方向上的性质可能不同。
\end{quote}

\subsubsection{1.3 晶体与非晶体的对比}\label{ux6676ux4f53ux4e0eux975eux6676ux4f53ux7684ux5bf9ux6bd4}

{\def\LTcaptype{none} % do not increment counter
\begin{longtable}[]{@{}lll@{}}
\toprule\noalign{}
特征 & 晶体 & 非晶体(玻璃态) \\
\midrule\noalign{}
\endhead
\bottomrule\noalign{}
\endlastfoot
微观结构 & 长程有序(周期性排列) & 短程有序、长程无序 \\
宏观外形 & 规则几何多面体 & 无规则外形 \\
熔点 & 固定(有明确的熔点) & 不固定(软化温度范围) \\
各向异性 & 显著 & 各向同性 \\
X射线衍射 & 有尖锐衍射峰 & 无尖锐衍射峰(弥散散射) \\
\end{longtable}
}

**F+V=E+2 (F:晶面 V:顶点 E:晶棱)

\begin{figure}[H]\centering
\includegraphics[width=0.55\textwidth,keepaspectratio]{puha-13-1-quartz-crystal.jpg}
\caption{石英晶体的不同外形与晶面}
\end{figure}

\begin{center}\rule{0.5\linewidth}{0.5pt}\end{center}

\subsection{二、点阵理论}\label{ux4e8cux70b9ux9635ux7406ux8bba}

\subsubsection{2.1 结构基元与点阵}\label{ux7ed3ux6784ux57faux5143ux4e0eux70b9ux9635}

晶体内部的质点在三维空间作周期性的有规则排列。为研究这种周期性，将晶体内原子分子等微粒在空间的周期排列的结果抽象为 \textbf{点阵(lattice)}。
点阵是一组无限的点，连结其中任意两点可得一向量，将各个点按此向量平移能使它复原。平移必须是按向量平行移动；点阵中每个点都具有完全相同的周围环境。

\textbf{结构基元}：晶体中周期性变化的具体内容。各个结构基元必须满足四个条件：化学组成相同、空间结构相同、排列取向相同、周围环境相同。

\textbf{示例}：在伸展的聚乙烯链 \(-\mathrm{CH}_2-\mathrm{CH}_2-\) 中，两个相邻的 \(-\mathrm{CH}_2-\) 虽组成和结构相同，但取向不同、周围环境不同，因此一个结构基元应为 \(-\mathrm{CH}_2-\mathrm{CH}_2-\)，而非单个 \(-\mathrm{CH}_2-\)。

\begin{figure}[H]\centering
\includegraphics[width=0.55\textwidth,keepaspectratio]{CH2.png}
\caption{聚乙烯长链高分子}
\end{figure}

\begin{quote}
\textbf{认知冲突}：「为什么聚乙烯的结构基元不是单个CH\textsubscript{2}？」---单个CH\textsubscript{2}虽然组成和结构完全相同，但相邻两个CH\textsubscript{2}的C---C键取向相反(一个''正放''、一个''倒放'')，周围环境也不相同(一个左边是CH\textsubscript{2}、右边是CH\textsubscript{2}，但取向不同)。所以必须取两个CH\textsubscript{2}作为一个结构基元。
\end{quote}

确定了结构基元后，用一个抽象的几何点(点阵点)代表每个结构基元。

\begin{quote}
\textbf{练一练：结构基元识别}
\end{quote}

{\def\LTcaptype{none} % do not increment counter
\begin{longtable}[]{@{}lll@{}}
\toprule\noalign{}
物质 & 结构基元是？ & 判断理由 \\
\midrule\noalign{}
\endhead
\bottomrule\noalign{}
\endlastfoot
NaCl晶体 & \_\_\_\_\_\_ & \_\_\_\_\_\_ \\
干冰(CO\textsubscript{2}(s)) & \_\_\_\_\_\_ & \_\_\_\_\_\_ \\
金属Cu & \_\_\_\_\_\_ & \_\_\_\_\_\_ \\
伸展的聚乙烯链 & \_\_\_\_\_\_ & \_\_\_\_\_\_ \\
CsCl晶体 & \_\_\_\_\_\_ & \_\_\_\_\_\_ \\
\end{longtable}
}

从晶体中无数结构单元中抽象出的几何点集合构成 \textbf{点阵(lattice)}。

\textbf{核心关系}：

\[
\text{晶体结构} = \text{点阵} + \text{结构基元}
\]

\begin{figure}[H]\centering
\includegraphics[width=0.55\textwidth,keepaspectratio]{图-点阵结构基元晶胞关系.png}
\caption{点阵·结构基元·晶胞三要素关系图}
\end{figure}

\textbf{点阵的基本特征}：点阵是按周期性规律在空间排布的一组无限多个点。按连接其中任意两点的向量进行平移时，能使点阵复原。点阵中每个点具有相同的周围环境。

\begin{figure}[H]\centering
\includegraphics[width=0.32\textwidth,height=0.24\textheight,keepaspectratio]{puha-13-2a-line-lattice.jpg}
\hspace{0.015\textwidth}
\includegraphics[width=0.32\textwidth,height=0.24\textheight,keepaspectratio]{puha-13-2b-plane-lattice.jpg}
\hspace{0.015\textwidth}
\includegraphics[width=0.32\textwidth,height=0.24\textheight,keepaspectratio]{spatial-lattice.jpg}
\end{figure}

\subsubsection{2.2 平面点阵与正当单位}\label{ux5e73ux9762ux70b9ux9635ux4e0eux6b63ux5f53ux5355ux4f4d}

平面点阵可划分为无数并置的平行四边形(平面格子)。划分规则：

{\def\LTcaptype{none} % do not increment counter
\begin{longtable}[]{@{}cc@{}}
\toprule\noalign{}
格子类型 & 每个平行四边形平均含点阵点数 \\
\midrule\noalign{}
\endhead
\bottomrule\noalign{}
\endlastfoot
素格子 & 1 个 \\
复格子 & \(>1\) 个 \\
\end{longtable}
}

\textbf{正当平面单位的选取标准}(按优先级)：
1. 与平面点阵对称性一致的平行四边形
2. 对称性尽可能高(直角尽可能多)
3. 包含点阵点数目尽可能少(面积尽可能小)

\begin{figure}[H]\centering
\includegraphics[width=0.55\textwidth,keepaspectratio]{two.png}
\caption{4种平面点阵的正当单位}
\end{figure}

\begin{quote}
\textbf{易错提醒}：「点阵点 = 原子位置？」---不是！点阵点是抽象的几何点，代表一个结构基元的位置。原子是实际物质。NaCl晶体中每个点阵点对应1个Na\textsuperscript{+}+1个Cl\textsuperscript{-}(1个NaCl)，而不是1个原子。点阵点之间可能是原子、离子或分子，具体由结构基元决定。
\end{quote}

\subsubsection{2.3 空间点阵与14种布拉维格子}\label{ux7a7aux95f4ux70b9ux9635ux4e0e14ux79cdux5e03ux62c9ux7ef4ux683cux5b50}

空间点阵可划分为无数并置的平行六面体(空间格子)。平行六面体单位中：

{\def\LTcaptype{none} % do not increment counter
\begin{longtable}[]{@{}cc@{}}
\toprule\noalign{}
点阵点位置 & 对每个单位的平均贡献 \\
\midrule\noalign{}
\endhead
\bottomrule\noalign{}
\endlastfoot
顶角 & \(1/8\) \\
棱上 & \(1/4\) \\
面上 & \(1/2\) \\
内部 & \(1\) \\
\end{longtable}
}

\textbf{正当空间单位的选取标准}：
1. 符合相应空间点阵的对称性
2. 棱与棱之间的直角数目尽量多
3. 体积尽量小

法国晶体学家布拉维(Bravais)证明：晶体的点阵共有 \textbf{7 种形状、14 种型式}，即 \textbf{14 种布拉维格子}。

{\def\LTcaptype{none} % do not increment counter
\begin{longtable}[]{@{}ccccc@{}}
\toprule\noalign{}
晶系 & 简单(P) & 体心(I) & 面心(F) & 底心(C) \\
\midrule\noalign{}
\endhead
\bottomrule\noalign{}
\endlastfoot
三斜 & 有 & --- & --- & --- \\
单斜 & 有 & --- & --- & 有 \\
正交 & 有 & 有 & 有 & 有 \\
四方 & 有 & 有 & --- & --- \\
六方 & 有 & --- & --- & --- \\
菱方(三方) & 有 & --- & --- & --- \\
立方 & 有 & 有 & 有 & --- \\
\end{longtable}
}

\begin{figure}[H]\centering
\includegraphics[width=0.32\textwidth,height=0.24\textheight,keepaspectratio]{16.png}
\hspace{0.015\textwidth}
\includegraphics[width=0.32\textwidth,height=0.24\textheight,keepaspectratio]{17.png}
\hspace{0.015\textwidth}
\includegraphics[width=0.32\textwidth,height=0.24\textheight,keepaspectratio]{18.png}
\end{figure}

\emph{图 6 14种布拉维格子}

\begin{quote}
\textbf{记忆要点}：并非每个晶系都有全部四种带心方式。例如立方晶系只有 P、I、F 三种(无底心)；四方晶系只有 P、I 两种(无底心和面心---因为四方面心可由更小的四方体心代替，四方底心可由更小的简单四方代替)。
\end{quote}

\begin{center}\rule{0.5\linewidth}{0.5pt}\end{center}

\subsection{三、七大晶系}\label{ux4e09ux4e03ux5927ux6676ux7cfb}

\subsubsection{3.1 晶胞参数}\label{ux6676ux80deux53c2ux6570}

\begin{figure}[H]\centering
\includegraphics[width=0.55\textwidth,keepaspectratio]{图-七大晶系判定决策树.png}
\caption{七大晶系判定决策树}
\end{figure}

晶胞的大小和形状由 \textbf{晶胞参数} 决定：三个棱长 \(a\)、\(b\)、\(c\)，以及棱间的夹角 \(\alpha\)(\(b\) 与 \(c\) 夹角)、\(\beta\)(\(a\) 与 \(c\) 夹角)、\(\gamma\)(\(a\) 与 \(b\) 夹角)。

\subsubsection{3.2 七大晶系的晶胞参数特征}\label{ux4e03ux5927ux6676ux7cfbux7684ux6676ux80deux53c2ux6570ux7279ux5f81}

{\def\LTcaptype{none} % do not increment counter
\begin{longtable}[]{@{}
  >{\centering\arraybackslash}p{(\linewidth - 8\tabcolsep) * \real{0.2273}}
  >{\raggedright\arraybackslash}p{(\linewidth - 8\tabcolsep) * \real{0.1818}}
  >{\raggedright\arraybackslash}p{(\linewidth - 8\tabcolsep) * \real{0.1818}}
  >{\centering\arraybackslash}p{(\linewidth - 8\tabcolsep) * \real{0.2273}}
  >{\raggedright\arraybackslash}p{(\linewidth - 8\tabcolsep) * \real{0.1818}}@{}}
\toprule\noalign{}
\begin{minipage}[b]{\linewidth}\centering
晶系
\end{minipage} & \begin{minipage}[b]{\linewidth}\raggedright
轴长关系
\end{minipage} & \begin{minipage}[b]{\linewidth}\raggedright
轴角关系
\end{minipage} & \begin{minipage}[b]{\linewidth}\centering
独立参数数
\end{minipage} & \begin{minipage}[b]{\linewidth}\raggedright
实例
\end{minipage} \\
\midrule\noalign{}
\endhead
\bottomrule\noalign{}
\endlastfoot
\textbf{立方} & \(a = b = c\) & \(\alpha = \beta = \gamma = 90^\circ\) & 1 & NaCl, CsCl, Cu, 金刚石 \\
\textbf{四方} & \(a = b \neq c\) & \(\alpha = \beta = \gamma = 90^\circ\) & 2 & TiO\textsubscript{2}(金红石), CaC\textsubscript{2} \\
\textbf{六方} & \(a = b \neq c\) & \(\alpha = \beta = 90^\circ,\ \gamma = 120^\circ\) & 2 & Mg, Zn, 石墨 \\
\textbf{菱方(三方)} & \(a = b = c\) & \(\alpha = \beta = \gamma < 120^\circ,\ \neq 90^\circ\) & 2 & CaCO\textsubscript{3}(方解石) \\
\textbf{正交} & \(a \neq b \neq c\) & \(\alpha = \beta = \gamma = 90^\circ\) & 3 & 斜方硫 \\
\textbf{单斜} & \(a \neq b \neq c\) & \(\alpha = \gamma = 90^\circ,\ \beta \neq 90^\circ\) & 4 & 石膏 \\
\textbf{三斜} & \(a \neq b \neq c\) & \(\alpha \neq \beta \neq \gamma\) & 6 & 蓝晶石 \\
\end{longtable}
}

晶胞体积的通式：

\[
V = abc\sqrt{1 - \cos^2\alpha - \cos^2\beta - \cos^2\gamma + 2\cos\alpha\cos\beta\cos\gamma}
\]

对于立方晶系：\(V = a^3\)；对于正交晶系：\(V = abc\)。

\subsubsection{3.3 晶系的对称性判定}\label{ux6676ux7cfbux7684ux5bf9ux79f0ux6027ux5224ux5b9a}

在竞赛中，判断晶系优先依据 \textbf{对称性} 而非仅看轴长关系。最高对称轴的类型决定了晶系归属：

{\def\LTcaptype{none} % do not increment counter
\begin{longtable}[]{@{}cc@{}}
\toprule\noalign{}
最高对称轴 & 晶系 \\
\midrule\noalign{}
\endhead
\bottomrule\noalign{}
\endlastfoot
四个 \(C_3\) 轴 & 立方 \\
一个 \(C_4\) 或 \(\overline{4}\) 轴 & 四方 \\
一个 \(C_6\) 或 \(\overline{6}\) 轴 & 六方 \\
一个 \(C_3\) 或 \(\overline{3}\) 轴 & 三方(菱方) \\
三个互相垂直的 \(C_2\) 轴或镜面 & 正交 \\
一个 \(C_2\) 轴或镜面 & 单斜 \\
无(只有 \(C_1\) 或 \(i\)) & 三斜 \\
\end{longtable}
}

\begin{center}\rule{0.5\linewidth}{0.5pt}\end{center}

\subsection{四、晶胞}\label{ux56dbux6676ux80de}

\subsubsection{4.1 晶胞的基本特征}\label{ux6676ux80deux7684ux57faux672cux7279ux5f81}

\textbf{晶胞(unit cell)} 是晶体结构中的基本重复单位。整个晶体按晶胞在三维空间周期性地重复排列，即无数晶胞 \textbf{无隙并置} 堆砌而成。

\begin{quote}
\textbf{注意}：晶胞必须是平行六面体。从晶体中取六方棱柱作为晶胞是错误的---六方晶系的正当晶胞仍为平行六面体(\(a = b \neq c,\ \gamma = 120^\circ\))。
\end{quote}

\subsubsection{4.2 素晶胞与复晶胞}\label{ux7d20ux6676ux80deux4e0eux590dux6676ux80de}

{\def\LTcaptype{none} % do not increment counter
\begin{longtable}[]{@{}cclcl@{}}
\toprule\noalign{}
类型 & 符号 & 特征 & 每个晶胞含点阵点数 & 示例 \\
\midrule\noalign{}
\endhead
\bottomrule\noalign{}
\endlastfoot
\textbf{素晶胞} & P & 晶体结构中最小的基本单元 & 1 & 简单立方(Po) \\
\textbf{体心晶胞} & I & 可作体心平移 \((1/2, 1/2, 1/2)\) & 2 & 金属钾(bcc) \\
\textbf{面心晶胞} & F & 可作面心平移 \((1/2, 1/2, 0)\) 等 & 4 & 金属铜(fcc) \\
\textbf{底心晶胞} & C & 可作底心平移 \((1/2, 1/2, 0)\) & 2 & 碘晶体(B底心) \\
\end{longtable}
}

\begin{figure}[H]\centering
\includegraphics[width=0.32\textwidth,height=0.24\textheight,keepaspectratio]{图-简单立方晶胞3D.png}
\hspace{0.015\textwidth}
\includegraphics[width=0.32\textwidth,height=0.24\textheight,keepaspectratio]{13-7b-simple-cubic-unit-cell-po.jpg}
\hspace{0.015\textwidth}
\includegraphics[width=0.32\textwidth,height=0.24\textheight,keepaspectratio]{图-体心立方晶胞3D.png}
\hspace{0.015\textwidth}
\includegraphics[width=0.32\textwidth,height=0.24\textheight,keepaspectratio]{图-面心立方晶胞3D.png}
\caption{素晶胞与复晶胞对比—左上为素晶胞P(8个顶角，1个原子)，右上为简单立方堆积(Po)，左下为体心立方I(2个原子)，右下为面心立方F(4个原子)}
\end{figure}

\textbf{体心平移检验}：将原晶胞框架的顶角平移到体心位置，若得到的新晶胞与原晶胞无差别，则为体心晶胞。

例如氯化铯(CsCl)虽然看起来有体心原子，但 Cs\textsuperscript{+} 与 Cl\textsuperscript{-} 环境不同，不能作体心平移 \箭 **不是** 体心晶胞，而是简单立方晶胞。

\subsubsection{4.3 晶胞中原子的分数坐标}\label{ux6676ux80deux4e2dux539fux5b50ux7684ux5206ux6570ux5750ux6807}

晶胞中原子的位置用 \textbf{分数坐标(fractional coordinates)} 表示：\((x, y, z)\)，表示向量 \(x\boldsymbol{a} + y\boldsymbol{b} + z\boldsymbol{c}\)。

\begin{figure}[H]\centering
\includegraphics[width=0.55\textwidth,keepaspectratio]{图-分数坐标三维可视化.png}
\caption{分数坐标三维可视化}
\end{figure}

\textbf{常见位置的分数坐标}：

{\def\LTcaptype{none} % do not increment counter
\begin{longtable}[]{@{}cc@{}}
\toprule\noalign{}
位置 & 坐标 \\
\midrule\noalign{}
\endhead
\bottomrule\noalign{}
\endlastfoot
顶角 & \((0, 0, 0)\) \\
体心 & \((1/2, 1/2, 1/2)\) \\
\(a\) 棱心 & \((1/2, 0, 0)\) \\
\(b\) 棱心 & \((0, 1/2, 0)\) \\
\(c\) 棱心 & \((0, 0, 1/2)\) \\
\(ab\) 面心 & \((1/2, 1/2, 0)\) \\
\(ac\) 面心 & \((1/2, 0, 1/2)\) \\
\(bc\) 面心 & \((0, 1/2, 1/2)\) \\
\end{longtable}
}

\begin{quote}
\textbf{注意}：分数坐标的取值区间为 \(0 \leqslant x, y, z < 1\)。由晶胞的平移性可知，坐标 \((1, 0, 0)\) 等价于 \((0, 0, 0)\)。
\end{quote}

\subsubsection{4.4 晶胞中的原子计数}\label{ux6676ux80deux4e2dux7684ux539fux5b50ux8ba1ux6570}

\begin{figure}[H]\centering
\includegraphics[width=0.55\textwidth,keepaspectratio]{图-晶胞原子计数3D.png}
\caption{晶胞原子计数—顶角(深灰)1/8、棱心(中灰)1/4、面心(浅灰)1/2、体内(白)1}
\end{figure}

{\def\LTcaptype{none} % do not increment counter
\begin{longtable}[]{@{}ccl@{}}
\toprule\noalign{}
原子位置 & 对晶胞的贡献 & 计数示例 \\
\midrule\noalign{}
\endhead
\bottomrule\noalign{}
\endlastfoot
顶角 & \(1/8\) & 立方晶胞：\(8 \times 1/8 = 1\) \\
棱心 & \(1/4\) & 立方晶胞：\(12 \times 1/4 = 3\) \\
面心 & \(1/2\) & 立方晶胞：\(6 \times 1/2 = 3\) \\
体内 & \(1\) & --- \\
\end{longtable}
}

\textbf{示例---金刚石晶胞}：包含 8 个碳原子(顶角 \(8 \times 1/8\) + 面心 \(6 \times 1/2\) + 体内 4 个)，坐标分别为：
- \((0, 0, 0),\ (1/2, 1/2, 0),\ (1/2, 0, 1/2),\ (0, 1/2, 1/2)\)
- \((3/4, 1/4, 1/4),\ (1/4, 3/4, 1/4),\ (1/4, 1/4, 3/4),\ (3/4, 3/4, 3/4)\)

\begin{figure}[H]\centering
\includegraphics[width=0.48\textwidth,height=0.32\textheight,keepaspectratio]{13-15c-diamond-unit-cell-atom-count.jpg}
\hspace{0.015\textwidth}
\includegraphics[width=0.48\textwidth,height=0.32\textheight,keepaspectratio]{13-15-diamond-crystal-structure-fcc-unit-cell.jpg}
\caption{金刚石的晶胞原子计数(左)与面心立方晶体结构(右)}
\end{figure}

\begin{center}\rule{0.5\linewidth}{0.5pt}\end{center}

\subsection{五、等同点系}\label{ux4e94ux7b49ux540cux70b9ux7cfb}

\subsubsection{5.1 定义}\label{ux5b9aux4e49}

晶体结构中，若一组点从其中各个点出发，沿同一指定方向观察到的物质环境和几何环境都相同，则这些点属于同一个 \textbf{等同点系(equivalent point system)}。

\subsubsection{5.2 拆分规则}\label{ux62c6ux5206ux89c4ux5219}

{\def\LTcaptype{none} % do not increment counter
\begin{longtable}[]{@{}cl@{}}
\toprule\noalign{}
步骤 & 操作 \\
\midrule\noalign{}
\endhead
\bottomrule\noalign{}
\endlastfoot
1. & 先把不同的原子相互拆开(不同元素肯定不属于同一等同点系) \\
2. & 相同的原子若周围环境不同，需拆分为多个等同点系 \\
3. & 拆分到各等同点系上原子的排列方式已相同时停止 \\
\end{longtable}
}

\subsubsection{5.3 应用示例}\label{ux5e94ux7528ux793aux4f8b}

\textbf{CsCl 晶体}：拆分为 2 个等同点系---Cl\textsuperscript{-} 按简单立方排列，Cs\textsuperscript{+} 按简单立方排列。每个等同点系上离子的排列方式即为点阵---简单立方点阵。从两个等同点系各取 1 个离子(1 个 Cl\textsuperscript{-} + 1 个 Cs\textsuperscript{+})= \textbf{结构基元}。

\textbf{金刚石晶体}：有两类环境不同的碳原子---晶胞顶点/面心处的 CI 和体内四面体空隙中的 CII。拆分为 2 个等同点系，各自按面心立方排列。点阵为 \textbf{面心立方点阵}。结构基元 = 1 个 CI + 1 个 CII = \textbf{2 个碳原子}。

\begin{figure}[H]\centering
\includegraphics[width=0.55\textwidth,keepaspectratio]{puha-13-15a-diamond-unit-cell.jpg}
\caption{金刚石晶胞—顶点/面心为CI(面心立方排列)，体内四面体空隙中为CII，两类C环境不同}
\end{figure}

\begin{figure}[H]\centering
\includegraphics[width=0.32\textwidth,height=0.24\textheight,keepaspectratio]{13-11a-nacl-cscl-zns-three-ab-structures.jpg}
\hspace{0.015\textwidth}
\includegraphics[width=0.32\textwidth,height=0.24\textheight,keepaspectratio]{13-11b-cscl-structure-unit-cell.jpg}
\hspace{0.015\textwidth}
\includegraphics[width=0.32\textwidth,height=0.24\textheight,keepaspectratio]{13-11-three-ab-ionic-crystal-structures-nacl-cscl-zns.jpg}
\caption{三种典型的 AB 型离子晶体结构(从左至右：NaCl 型、CsCl 型、立方 ZnS 型)}
\end{figure}

\subsubsection{5.4 等同点系与点阵的区别}\label{ux7b49ux540cux70b9ux7cfbux4e0eux70b9ux9635ux7684ux533aux522b}

\begin{quote}
\textbf{易错提醒}：不能认为点阵是从等同点系抽象出来的。任何晶体都只能抽象出 \textbf{唯一确定} 的点阵，但可能有多个等同点系。一个点阵点代表的结构基元必须从每个等同点系上各取 1 个原子集合而成。
\end{quote}

\begin{quote}
\textbf{认知冲突}：「等同点系与点阵」---这是晶体学中最易混淆的概念之一。\textbf{点阵}是抽象几何点，\textbf{等同点系}是晶体中实际原子位置的集合。点阵是''骨架''(几何位置关系)，等同点系是''骨架上的内容''。不同等同点系可以有相同的排列方式(如CsCl中Cl\textsuperscript{-}和Cs\textsuperscript{+}都按简单立方排列)，但它们属于不同的等同点系。
\end{quote}

\begin{center}\rule{0.5\linewidth}{0.5pt}\end{center}

\subsection{六、对称元素基础}\label{ux516dux5bf9ux79f0ux5143ux7d20ux57faux7840}

\subsubsection{6.1 晶体中的对称元素}\label{ux6676ux4f53ux4e2dux7684ux5bf9ux79f0ux5143ux7d20}

{\def\LTcaptype{none} % do not increment counter
\begin{longtable}[]{@{}
  >{\centering\arraybackslash}p{(\linewidth - 6\tabcolsep) * \real{0.0870}}
  >{\centering\arraybackslash}p{(\linewidth - 6\tabcolsep) * \real{0.1522}}
  >{\raggedright\arraybackslash}p{(\linewidth - 6\tabcolsep) * \real{0.2609}}
  >{\raggedright\arraybackslash}p{(\linewidth - 6\tabcolsep) * \real{0.5000}}@{}}
\toprule\noalign{}
\begin{minipage}[b]{\linewidth}\centering
对称元素
\end{minipage} & \begin{minipage}[b]{\linewidth}\centering
符号
\end{minipage} & \begin{minipage}[b]{\linewidth}\raggedright
操作
\end{minipage} & \begin{minipage}[b]{\linewidth}\raggedright
说明
\end{minipage} \\
\midrule\noalign{}
\endhead
\bottomrule\noalign{}
\endlastfoot
\textbf{旋转轴} & \(C_n\) & 绕轴旋转 \(360^\circ/n\) 复原 & \(n = 1, 2, 3, 4, 6\)(晶体中无 \(C_5\) 及 \(n > 6\) 的旋转轴) \\
\textbf{镜面} & \(\sigma\) & 通过镜面反映 & 镜面将晶体分为互为镜像的两部分 \\
\textbf{对称中心} & \(i\) & 通过中心反演 & 将 \((x, y, z)\) 变换为 \((-x, -y, -z)\) \\
\textbf{映轴} & \(S_n\) & 旋转 \(360^\circ/n\) 后对垂直面反映 & \(S_4\) 在晶体中出现(如 CH\textsubscript{4} 的正四面体对称性) \\
\textbf{反轴} & \(\overline{n}\) & 旋转 \(360^\circ/n\) 后反演 & \(\overline{1} = i\)，\(\overline{2} = \sigma\)(镜面) \\
\end{longtable}
}

\begin{figure}[H]\centering
\includegraphics[width=0.55\textwidth,keepaspectratio]{n.png}
\caption{反轴 $\overline{4}$}
\end{figure}

\begin{figure}[H]\centering
\includegraphics[width=0.55\textwidth,keepaspectratio]{atkins-cube-symmetry-elements.jpg}
\caption{立方体中的对称元素—二重轴(椭圆标记)、三重轴(三角标记)、四重轴(正方形标记)(Atkins 物理化学 图10A.1)}
\end{figure}

\begin{figure}[H]\centering
\includegraphics[width=0.55\textwidth,keepaspectratio]{atkins-nh3-h2o-symmetry.jpg}
\caption{分子中的对称元素示例—(a) NH 具有一个三重轴 $C_3$；(b) HO 具有一个二重轴 $C_2$(Atkins 物理化学 图10A.2)}
\end{figure}

\subsubsection{6.2 晶体的对称性限制}\label{ux6676ux4f53ux7684ux5bf9ux79f0ux6027ux9650ux5236}

\begin{figure}[H]\centering
\includegraphics[width=0.55\textwidth,keepaspectratio]{图-对称性限制定律.png}
\caption{晶体对称性限制定律}
\end{figure}

晶体中旋转轴的轴次只能为 \(n = 1, 2, 3, 4, 6\)，不可能出现 \(C_5\) 及 \(n > 6\) 的旋转轴。这源于平移对称性的约束---称为 \textbf{晶体对称性限制定律}。

\begin{quote}
\textbf{认知冲突}：「准晶 (Quasicrystal) 打破了对称性限制定律？」---不完全是！1982年 Shechtman 在急冷 Al-Mn 合金中发现了5次对称衍射图案，这似乎''打破''了传统晶体学定律。但准晶并非周期性晶体---它是有序但非周期性的结构(利用彭罗斯铺砌的数学原理)。这一发现迫使国际晶体学联合会重新定义''晶体''，Shechtman 也因此获得2011年诺贝尔化学奖。
\end{quote}

\subsubsection{6.3 晶系与对称性的对应}\label{ux6676ux7cfbux4e0eux5bf9ux79f0ux6027ux7684ux5bf9ux5e94}

晶体的七个晶系由其最高对称轴唯一确定(见 §3.3 的表)。竞赛中常见的对应关系：

{\def\LTcaptype{none} % do not increment counter
\begin{longtable}[]{@{}lcl@{}}
\toprule\noalign{}
对称性特征 & 晶系 & 典型实例 \\
\midrule\noalign{}
\endhead
\bottomrule\noalign{}
\endlastfoot
有 4 个 \(C_3\) 轴 & 立方 & 金刚石、NaCl、Cu \\
有 1 个 \(C_4\) 轴 & 四方 & TiO\textsubscript{2}(金红石)、SnO\textsubscript{2} \\
有 1 个 \(C_6\) 轴 & 六方 & 石墨、Zn(hcp) \\
有 1 个 \(C_3\) 轴 & 三方 & CaCO\textsubscript{3}(方解石) \\
\end{longtable}
}

\begin{quote}
\textbf{深化}：晶体的完整对称性分类(32 种点群、230 种空间群)属于结构化学深化内容。
\end{quote}

\begin{center}\rule{0.5\linewidth}{0.5pt}\end{center}

\subsection{七、规律与方法}\label{ux4e03ux89c4ux5f8bux4e0eux65b9ux6cd5}

\subsubsection{7.1 晶体学分析的三步流程}\label{ux6676ux4f53ux5b66ux5206ux6790ux7684ux4e09ux6b65ux6d41ux7a0b}

\begin{verbatim}
宏观晶体外形 → 对称性判断 → 晶系归属
    ↓
抽象点阵 → 点阵型式（P/I/F/C）→ 布拉维格子
    ↓
确定晶胞 → 原子坐标 → 原子计数 → 结构基元
\end{verbatim}

\subsubsection{7.2 常见误区}\label{ux5e38ux89c1ux8befux533a}

\begin{quote}
\textbf{误区 1}：``晶胞就是最小的结构单元''---晶胞是点阵的基本重复单位，但不一定是最小的(素晶胞才是最小的)。复晶胞(I、F、C)是素晶胞的多倍体。

\textbf{误区 2}：``CsCl 是体心立方''---CsCl 看似体心立方，但 Cs\textsuperscript{+} 与 Cl\textsuperscript{-} 环境不同，不能作体心平移，应归为简单立方点阵。

\textbf{误区 3}：``不存在有 \(C_5\) 轴的分子''---分子可以具有 \(C_5\) 轴(如二茂铁)，但晶体中不允许出现 \(C_5\) 轴。具有 \(C_5\) 轴的分子在结晶时，其分子对称性与晶体的平移对称性不兼容。

\textbf{误区 4}：``所有晶系都有 14 种布拉维格子''---14 种是所有晶系的格子总和。每个晶系只有特定几种带心方式(如立方只有 P/I/F，四方只有 P/I)。
\end{quote}

\begin{center}\rule{0.5\linewidth}{0.5pt}\end{center}

\subsection{八、典型例题}\label{ux516bux5178ux578bux4f8bux9898}

\textbf{题目}：某晶体的晶胞参数为 \(a = 0.520\ \mathrm{nm},\ b = 0.520\ \mathrm{nm},\ c = 0.856\ \mathrm{nm},\ \alpha = \beta = \gamma = 90^\circ\)。判断该晶体属于什么晶系，有多少个独立晶胞参数。

(答案见§十二)

\begin{center}\rule{0.5\linewidth}{0.5pt}\end{center}

\textbf{题目}：金属钠的晶体结构为体心立方(bcc)，每个晶胞含多少个钠原子？写出体心立方晶胞中所有原子的分数坐标。

(答案见§十二)

\begin{center}\rule{0.5\linewidth}{0.5pt}\end{center}

\textbf{题目}：CaC\textsubscript{2} 晶体结构如下图所示，其中 C---C 键长为 0.120 nm，C\textsubscript{2}\textsuperscript{2-} 离子呈哑铃型且键轴平行于 c 轴排列。问 CaC\textsubscript{2} 的正当晶胞是否仍为面心立方？属于什么点阵型式？

\begin{figure}[H]\centering
\includegraphics[width=0.55\textwidth,keepaspectratio]{cac2-unit-cell.png}
\end{figure}

\textbf{思路分析}：
1. CaC\textsubscript{2} 初看类似 NaCl 型结构，Na\textsuperscript{+} 替换为 Ca\textsuperscript{2+}，Cl\textsuperscript{-} 替换为 C\textsubscript{2}\textsuperscript{2-}
2. 但 C\textsubscript{2}\textsuperscript{2-} 为哑铃型，其键轴平行于 c 轴 \箭 破坏了立方对称性
3. 需要从对称性角度判断晶系归属和点阵型式

(答案见§十二)

\begin{center}\rule{0.5\linewidth}{0.5pt}\end{center}

\textbf{题目}：试分析 Cu\textsubscript{2}O(赤铜矿)的晶胞是否为体心立方？已知 Cu\textsubscript{2}O 晶胞中 O 原子坐标为 \((0, 0, 0)\) 和 \((1/2, 1/2, 1/2)\)；Cu 原子坐标为 \((3/4, 1/4, 1/4)\)、\((1/4, 3/4, 1/4)\)、\((1/4, 1/4, 3/4)\)、\((3/4, 3/4, 3/4)\)。

(答案见§十二)

\begin{center}\rule{0.5\linewidth}{0.5pt}\end{center}

\textbf{题目}：某晶体的粉末衍射数据显示其属于立方晶系，晶胞参数 \(a = 543\ \mathrm{pm}\)。已知该晶体中原子排列的对称性显示其具有面心平移特征(F)，判读其每个晶胞含 4 个原子。若该晶体的密度为 \(2.33\ \mathrm{g\cdot cm^{-3}}\)，求其摩尔质量，并推断可能的物质。

(答案见§十二)

\begin{center}\rule{0.5\linewidth}{0.5pt}\end{center}

\subsection{十、本节总结 · 速查}\label{ux5341ux672cux8282ux603bux7ed3-ux901fux67e5}

{\def\LTcaptype{none} % do not increment counter
\begin{longtable}[]{@{}
  >{\raggedright\arraybackslash}p{(\linewidth - 4\tabcolsep) * \real{0.1548}}
  >{\raggedright\arraybackslash}p{(\linewidth - 4\tabcolsep) * \real{0.4643}}
  >{\raggedright\arraybackslash}p{(\linewidth - 4\tabcolsep) * \real{0.3810}}@{}}
\toprule\noalign{}
\begin{minipage}[b]{\linewidth}\raggedright
核心概念
\end{minipage} & \begin{minipage}[b]{\linewidth}\raggedright
关键规律
\end{minipage} & \begin{minipage}[b]{\linewidth}\raggedright
常见陷阱
\end{minipage} \\
\midrule\noalign{}
\endhead
\bottomrule\noalign{}
\endlastfoot
\textbf{晶体定义} & 长程有序 + 周期性排列 & 非晶体也有短程有序但长程无序 \\
\textbf{晶体性质} & 自范性、各向异性、对称性、均一性 & 各向异性 \不等于 不均一 \\
\textbf{点阵} & 晶体结构 = 点阵 + 结构基元 & 等同点系 \不等于 点阵(晶体只有一个点阵) \\
\textbf{14 种布拉维格子} & 7 种晶系 × 带心方式(P/I/F/C) & 并非每种晶系都有全部 4 种 \\
\textbf{晶胞计数} & 顶角 \(1/8\)、棱心 \(1/4\)、面心 \(1/2\)、体内 \(1\) & CsCl 不是体心(体心平移条件不满足) \\
\textbf{七大晶系} & 由最高对称轴决定 & 仅凭轴长轴角判断有时不可靠 \\
\textbf{对称元素} & \(C_n\)、\(\sigma\)、\(i\)、\(S_n\)、\(\overline{n}\) & 晶体中无 \(C_5\)(对称性限制定律) \\
\textbf{素/复晶胞} & P = 1 个点阵点，I = 2，F = 4，C = 2 & 分数坐标 \((1, 0, 0)\) 等价于 \((0, 0, 0)\) \\
\end{longtable}
}

\begin{center}\rule{0.5\linewidth}{0.5pt}\end{center}

\subsection{十一、练习题}\label{ux5341ux4e00ux7ec3ux4e60ux9898}

\begin{enumerate}
\def\labelenumi{\arabic{enumi}.}
\item
  \textbf{晶体性质}：请解释为什么晶体具有固定的熔点而非晶体没有。
\item
  \textbf{七大晶系判断}：根据以下晶胞参数判断所属晶系：

  \begin{enumerate}
  \def\labelenumii{(\arabic{enumii})}
  \tightlist
  \item
    \(a = b = c = 0.400\ \mathrm{nm},\ \alpha = \beta = \gamma = 90^\circ\)
  \item
    \(a = b = 0.500\ \mathrm{nm},\ c = 0.800\ \mathrm{nm},\ \alpha = \beta = 90^\circ,\ \gamma = 120^\circ\)
  \item
    \(a = 0.700\ \mathrm{nm},\ b = 0.700\ \mathrm{nm},\ c = 0.900\ \mathrm{nm},\ \alpha = \beta = \gamma = 90^\circ\)
  \end{enumerate}
\item
  \textbf{晶胞计数}：计算面心立方(fcc)晶胞中有多少个原子。写出 fcc 晶胞中所有原子的分数坐标。
\item
  \textbf{点阵型式判断}：某化合物的晶胞中，A 原子位于顶点，B 原子位于面心，C 原子位于体心。判断：

  \begin{enumerate}
  \def\labelenumii{(\arabic{enumii})}
  \tightlist
  \item
    该晶胞属于简单、体心、面心还是底心？
  \item
    写出该晶体的化学式。
  \item
    写出所有原子的分数坐标。
  \end{enumerate}
\item
  \textbf{等同点系应用}：已知金刚石晶胞中碳原子坐标为 \((0, 0, 0)\)、\((1/2, 1/2, 0)\)、\((1/2, 0, 1/2)\)、\((0, 1/2, 1/2)\)、\((3/4, 1/4, 1/4)\)、\((1/4, 3/4, 1/4)\)、\((1/4, 1/4, 3/4)\)、\((3/4, 3/4, 3/4)\)。试证明金刚石晶体的点阵为面心立方，结构基元含 2 个碳原子。
\item
  \textbf{对称性与晶系}：某晶体的晶胞有三个互相垂直的棱，长度分别为 \(a = 0.510\ \mathrm{nm}\)、\(b = 0.430\ \mathrm{nm}\)、\(c = 0.640\ \mathrm{nm}\)，棱间夹角均为 \(90^\circ\)。该晶体有几个对称元素？属于什么晶系？
\item
  \textbf{综合推理}：金属钨为体心立方(bcc)结构，晶胞参数 \(a = 0.316\ \mathrm{nm}\)，密度 \(\rho = 19.3\ \mathrm{g\cdot cm^{-3}}\)。

  \begin{enumerate}
  \def\labelenumii{(\arabic{enumii})}
  \tightlist
  \item
    计算每个晶胞中的钨原子数。
  \item
    计算阿伏加德罗常数 \(N_A\)。
  \item
    计算空间利用率(堆积效率)。
  \end{enumerate}
\item
  \textbf{晶系归属判定}：已知 CaC\textsubscript{2} 的晶体结构在室温下为体心四方(I)，晶胞参数 \(a = 0.387\ \mathrm{nm}\)、\(c = 0.637\ \mathrm{nm}\)。当温度升高到 \(450^\circ\mathrm{C}\) 以上时，CaC\textsubscript{2} 转变为面心立方结构，晶胞参数 \(a = 0.580\ \mathrm{nm}\)。试解释：

  \begin{enumerate}
  \def\labelenumii{(\arabic{enumii})}
  \tightlist
  \item
    为什么室温时 CaC\textsubscript{2} 为四方晶系而非立方晶系？
  \item
    高温时对称性升高为立方晶系，C\textsubscript{2}\textsuperscript{2-} 离子取向可能发生什么变化？
  \item
    计算两种晶型下 CaC\textsubscript{2} 的理论密度，比较大小。
  \end{enumerate}
\end{enumerate}

\begin{center}\rule{0.5\linewidth}{0.5pt}\end{center}

\subsection{十二、参考答案与提示}\label{ux5341ux4e8cux53c2ux8003ux7b54ux6848ux4e0eux63d0ux793a}

\textbf{11.1 基础巩固}

\begin{enumerate}
\def\labelenumi{\arabic{enumi}.}
\item
  晶体具有长程有序的周期性结构，所有原子/离子处于相同类型的环境中，加热时各处的原子同时获得足够能量脱离晶格位置 \箭 在固定温度全部熔化。非晶体无长程有序，各处结合能不同 \箭 随温度升高逐渐软化，无固定熔点。
\item
  \begin{enumerate}
  \def\labelenumii{(\arabic{enumii})}
  \tightlist
  \item
    立方晶系；(2) 六方晶系；(3) 四方晶系。
  \end{enumerate}
\item
  fcc 晶胞含 \(8 \times 1/8 + 6 \times 1/2 = 4\) 个原子。分数坐标：\((0, 0, 0)\)、\((1/2, 1/2, 0)\)、\((1/2, 0, 1/2)\)、\((0, 1/2, 1/2)\)。
\end{enumerate}

\textbf{11.2 提高训练}

\begin{enumerate}
\def\labelenumi{\arabic{enumi}.}
\setcounter{enumi}{3}
\item
  \textbf{提示}：(1) A 在顶点、B 在面心、C 在体心 \箭 顶点平移至体心后结构不变(A 仍到 A，B 仍到 B)\箭 体心晶胞(I)。(2) 化学式：\(A_{8\times 1/8} B_{6\times 1/2} C_{1\times 1} = \mathrm{AB_3C}\)。(3) A: \((0,0,0)\)；B: \((1/2,1/2,0)\)、\((1/2,0,1/2)\)、\((0,1/2,1/2)\)；C: \((1/2,1/2,1/2)\)。
\item
  \textbf{提示}：前 4 个原子(CI)坐标位于顶点和面心 \箭 面心立方排列。后 4 个原子(CII)坐标相对于 CI偏移 \((1/4, 1/4, 1/4)\) \箭 亦为面心立方排列。两类碳原子环境不同，分属 2 个等同点系 \箭 点阵为面心立方。每个面心立方点阵点对应 1 个 CI + 1 个 CII = 2 个 C 原子 \箭 结构基元含 2 个碳原子。
\item
  \textbf{提示}：\(a \neq b \neq c\)，\(\alpha = \beta = \gamma = 90^\circ\) \箭 正交晶系。正交晶系在三个互相垂直的方向上各有 \(C_2\) 轴(或镜面)。
\end{enumerate}

\textbf{11.3 挑战题}

\begin{enumerate}
\def\labelenumi{\arabic{enumi}.}
\setcounter{enumi}{6}
\item
  \textbf{提示}：(1) bcc 晶胞含 2 个钨原子。(2) 由 \(\rho = \dfrac{2M}{N_A a^3}\)，代入 \(M_{\mathrm{W}} = 183.84\ \mathrm{g\cdot mol^{-1}}\)、\(a = 0.316\ \mathrm{nm} = 3.16\times 10^{-8}\ \mathrm{cm}\)，可解得 \(N_A\)。(3) bcc 空间利用率 = 68\%。
\item
  \textbf{提示}：(1) 室温时 C\textsubscript{2}\textsuperscript{2-} 的 C---C 键轴平行于 c 轴排列 \箭 破坏了立方晶系的四个 \(C_3\) 轴对称性，保留 \(C_4\) 轴 \箭 四方晶系。(2) 高温时 C\textsubscript{2}\textsuperscript{2-} 热运动加剧，C---C 键轴取向无序化(在各方向上随机分布)\箭 宏观上呈现出立方对称性。(3) 计算密度 \(\rho = nM/N_A V\)，注意两种晶型晶胞所含 CaC\textsubscript{2} 的式量数不同(面心立方时 \(Z = 4\)，体心四方时 \(Z = 2\))，对比结果可知密度变化。
\end{enumerate}

\textbf{例题答案(§八)}

\begin{enumerate}
\def\labelenumi{\arabic{enumi}.}
\item
  \(a = b \neq c\)，\(\alpha = \beta = \gamma = 90^\circ\) \箭 **四方晶系\textbf{；独立参数 }2 个**。
\item
  bcc 晶胞含 \textbf{2 个} 原子。分数坐标：\((0, 0, 0)\) 和 \((1/2, 1/2, 1/2)\)。
\item
  不是面心立方；C\textsubscript{2}\textsuperscript{2-} 键轴平行于 c 轴 \箭 四方晶系；正当晶胞为 \textbf{体心四方(I)}。
\item
  进行体心平移检验：O 原子位置不变，但 Cu 原子位置改变 \箭 **不具备体心平移特征** \箭 不是体心立方晶胞。核心标准：体心平移后晶体结构是否不变。
\item
  \(M \approx 28.1\ \mathrm{g\cdot mol^{-1}}\) \箭 **硅(Si)**。
\end{enumerate}

\begin{center}\rule{0.5\linewidth}{0.5pt}\end{center}

\begin{quote}
\textbf{过渡}：上篇建立了晶体学的概念语言(点阵、晶胞、晶系、对称性)，下篇将用这些工具分析具体的晶体结构---金属密堆积、离子晶体结构类型、半径比规则与Pauling规则。
\end{quote}

\begin{center}\rule{0.5\linewidth}{0.5pt}\end{center}

\subsection{十四、金属密堆积}\label{ux5341ux56dbux91d1ux5c5eux5bc6ux5806ux79ef}

\begin{quote}
\textbf{关联知识点}：密堆积 · 金属晶体 · 配位数
\end{quote}

\subsubsection{14.1 为什么金属原子要''挤''在一起？}\label{ux4e3aux4ec0ux4e48ux91d1ux5c5eux539fux5b50ux8981ux6324ux5728ux4e00ux8d77}

金属原子的核外电子分布决定了它们倾向于以\textbf{尽可能密集}的方式排列---因为密集排列时原子间的距离最小，金属键(自由电子气与正离子实之间的静电吸引)最强，体系能量最低。所以绝大多数金属晶体的结构，就是\textbf{等径圆球的最密堆积}。

\subsubsection{14.2 四种基本堆积方式}\label{ux56dbux79cdux57faux672cux5806ux79efux65b9ux5f0f}

等径圆球在平面上排列时，有\textbf{密置层}(最紧密排列，每个球周围6个球)和\textbf{非密置层}(每个球周围4个球，空隙较大)两种基本方式。三维堆积就是在这些二维层的基础上叠加而成。

\begin{center}
\begin{minipage}[t]{0.48\textwidth}\centering
\includegraphics[width=0.95\linewidth,height=0.30\textheight,keepaspectratio]{puha-13-6a-non-close-packed-layer.jpg}
\captionof{figure}{非密置层—球心呈正方形排列，空隙较大}
\end{minipage}
\hfill
\begin{minipage}[t]{0.48\textwidth}\centering
\includegraphics[width=0.95\linewidth,height=0.30\textheight,keepaspectratio]{puha-13-6b-close-packed-layer.jpg}
\captionof{figure}{密置层—球心呈正三角形排列，每个球与6个球相切}
\end{minipage}
\end{center}

\begin{center}
\begin{minipage}[t]{0.48\textwidth}\centering
\includegraphics[width=0.95\linewidth,height=0.30\textheight,keepaspectratio]{puha-13-7b-sc-unit-cell.jpg}
\captionof{figure}{简单立方晶胞—非密置层逐层正对叠加(配位数6，空间利用率52%)}
\end{minipage}
\hfill
\begin{minipage}[t]{0.48\textwidth}\centering
\includegraphics[width=0.95\linewidth,height=0.30\textheight,keepaspectratio]{puha-13-7b-bcc-unit-cell.jpg}
\captionof{figure}{体心立方晶胞—非密置层逐层错位叠加(配位数8，空间利用率68%)}
\end{minipage}
\end{center}

{\def\LTcaptype{none} % do not increment counter
\begin{longtable}[]{@{}cccccl@{}}
\toprule\noalign{}
堆积类型 & 符号 & 配位数 & 晶胞原子数 & 空间利用率 & 实例 \\
\midrule\noalign{}
\endhead
\bottomrule\noalign{}
\endlastfoot
简单立方 & sc & 6 & 1 & 52.36\% & Po(钋) \\
体心立方 & bcc / A2 & 8 & 2 & 68.02\% & Fe、Na、K、Cr \\
面心立方 & fcc / ccp / A1 & 12 & 4 & 74.05\% & Cu、Ag、Au、Al \\
六方密堆积 & hcp / A3 & 12 & 2 & 74.05\% & Mg、Zn、Ti、Co \\
\end{longtable}
}

\begin{center}
\begin{minipage}[t]{0.48\textwidth}\centering
\includegraphics[width=0.95\linewidth,height=0.30\textheight,keepaspectratio]{ccp-unit-cell.jpg}
\captionof{figure}{面心立方晶胞(fcc/ccp)}
\end{minipage}
\hfill
\begin{minipage}[t]{0.48\textwidth}\centering
\includegraphics[width=0.95\linewidth,height=0.30\textheight,keepaspectratio]{hcp-unit-cell.jpg}
\captionof{figure}{六方晶胞(hcp)}
\end{minipage}
\end{center}

\begin{quote}
\textbf{认知冲突}：「为什么ccp和hcp的堆积系数都是74.05\%，但晶系不同？」---两者的差异不在堆积密度，而在\textbf{堆积层序}：ccp是ABCABC\ldots(三层一循环，产生立方对称性)，hcp是ABAB\ldots(两层一循环，六方对称性)。\textbf{同样的密度，不同的排列方式}。
\end{quote}

\begin{quote}
\textbf{课堂原话}：``CCP=ABC, HCP=AB, 密度一样高，堆法各不同。第三层去了新位置是CCP(立方)，回到原位是HCP(六方)。''
\end{quote}

\begin{center}
\begin{minipage}[t]{0.48\textwidth}\centering
\includegraphics[width=0.95\linewidth,height=0.30\textheight,keepaspectratio]{puha-13-9b-ccp-layer-relation.jpg}
\captionof{figure}{CCP层序—第三层落在C位(ABCABC…)}
\end{minipage}
\hfill
\begin{minipage}[t]{0.48\textwidth}\centering
\includegraphics[width=0.95\linewidth,height=0.30\textheight,keepaspectratio]{puha-13-9b-hcp-layer-relation.jpg}
\captionof{figure}{HCP层序—第三层回到A位(ABAB…)}
\end{minipage}
\end{center}

\begin{center}
\begin{minipage}[t]{0.48\textwidth}\centering
\includegraphics[width=0.95\linewidth,height=0.30\textheight,keepaspectratio]{puha-13-9c-fcc-unit-cell-from-packing.jpg}
\captionof{figure}{从CCP堆积中取出面心立方晶胞}
\end{minipage}
\hfill
\begin{minipage}[t]{0.48\textwidth}\centering
\includegraphics[width=0.95\linewidth,height=0.30\textheight,keepaspectratio]{puha-13-9c-hex-unit-cell-from-packing.jpg}
\captionof{figure}{从HCP堆积中取出六方晶胞}
\end{minipage}
\end{center}

\subsubsection{14.3 空间利用率全推导}\label{ux7a7aux95f4ux5229ux7528ux7387ux5168ux63a8ux5bfc}

\begin{quote}
\textbf{关键思路}：空间利用率 = \(\frac{\text{晶胞中所有原子体积之和}}{\text{晶胞体积}} = \frac{Z \times \frac{4}{3}\pi r^3}{a^3}\)
核心步骤：1. 找出原子接触方向 2. 用原子半径 r 表示晶胞参数 a 3. 代入公式计算
\end{quote}

\textbf{14.3.1 面心立方(fcc/ccp)}

\textbf{原子接触方向}：面对角线方向(因为面心原子与顶点原子相切)

\(\sqrt{2}a = 4r\) \箭 \(a = 2\sqrt{2}r\)

\textbf{晶胞原子数}：顶点 \(8 \times 1/8 = 1\) + 面心 \(6 \times 1/2 = 3\) = \textbf{4个}

\[
\text{空间利用率} = \frac{4 \times \frac{4}{3}\pi r^3}{(2\sqrt{2}r)^3} = \frac{16\pi r^3 / 3}{16\sqrt{2}r^3} = \frac{\pi}{3\sqrt{2}} \approx 74.05\%
\]

\textbf{14.3.2 体心立方(bcc)}

\textbf{原子接触方向}：体对角线方向(体心原子与顶点原子相切)

\(\sqrt{3}a = 4r\) \箭 \(a = \frac{4r}{\sqrt{3}}\)

\textbf{晶胞原子数}：顶点 \(8 \times 1/8 = 1\) + 体心 \(1\) = \textbf{2个}

\[
\text{空间利用率} = \frac{2 \times \frac{4}{3}\pi r^3}{(4r/\sqrt{3})^3} = \frac{8\pi r^3 / 3}{64r^3 / 3\sqrt{3}} = \frac{\sqrt{3}\pi}{8} \approx 68.02\%
\]

\textbf{14.3.3 六方密堆积(hcp)}

\textbf{原子接触方向}：a轴方向(底面正六边形边长 = 2r)，c轴方向的层间距决定理想轴比

底面正六边形面积：\(S = \frac{3\sqrt{3}}{2}a^2 = \frac{3\sqrt{3}}{2}(2r)^2 = 6\sqrt{3}r^2\)

理想轴比(两层之间的垂直距离 = 两个正四面体的高)：\(c = 2 \times \sqrt{\frac{2}{3}}a = 4\sqrt{\frac{2}{3}}r\)

\[
V = S \times c = 6\sqrt{3}r^2 \times 4\sqrt{\frac{2}{3}}r = 24\sqrt{2}r^3
\]

\textbf{晶胞原子数}：12个顶点各1/6(简单六方中)\箭 晶胞含2个原子

\[
\text{空间利用率} = \frac{2 \times \frac{4}{3}\pi r^3}{24\sqrt{2}r^3} = \frac{\pi}{3\sqrt{2}} \approx 74.05\%
\]

\begin{quote}
\textbf{练一练}：推导hcp的理想轴比 \(c/a\)。提示：密置双层中，上层球正好落在下层三个球围成的凹陷处，构成正四面体。两层球的球心垂直距离 = 正四面体的高 = \(\sqrt{6}/3 \times\) 边长。
(答案见§二十二参考答案)
\end{quote}

\subsubsection{14.4 空隙分析}\label{ux7a7aux9699ux5206ux6790}

金属堆积中，原子之间不可避免会留下空隙(即使是最密堆积)。这些空隙是理解离子晶体结构的钥匙---\textbf{正离子就是填充在负离子堆积形成的空隙中}。

{\def\LTcaptype{none} % do not increment counter
\begin{longtable}[]{@{}
  >{\centering\arraybackslash}p{(\linewidth - 6\tabcolsep) * \real{0.2632}}
  >{\centering\arraybackslash}p{(\linewidth - 6\tabcolsep) * \real{0.2632}}
  >{\centering\arraybackslash}p{(\linewidth - 6\tabcolsep) * \real{0.2632}}
  >{\raggedright\arraybackslash}p{(\linewidth - 6\tabcolsep) * \real{0.2105}}@{}}
\toprule\noalign{}
\begin{minipage}[b]{\linewidth}\centering
堆积类型
\end{minipage} & \begin{minipage}[b]{\linewidth}\centering
四面体空隙数
\end{minipage} & \begin{minipage}[b]{\linewidth}\centering
八面体空隙数
\end{minipage} & \begin{minipage}[b]{\linewidth}\raggedright
说明
\end{minipage} \\
\midrule\noalign{}
\endhead
\bottomrule\noalign{}
\endlastfoot
\textbf{fcc/ccp} & 8(每个晶胞) & 4(每个晶胞) & 每个球对应 1 个八面体空隙 + 2 个四面体空隙 \\
\textbf{hcp} & 每个球对应 2 个 & 每个球对应 1 个 & 与 fcc 空隙数相同，但排列方式不同 \\
\textbf{bcc} & 无正四面体空隙 & 无正八面体空隙 & 空隙形状为非正多面体(变形的四面体/八面体) \\
\end{longtable}
}

\textbf{空隙数快速验证}：fcc 中，每个球周围有 8 个四面体空隙(4个在上方+4个在下方)，但每个四面体空隙由 4 个球共享：\(8 \times 4 / 4 = 8\) 个四面体空隙。每个球周围有 6 个八面体空隙，每个由 6 个球共享：\(6 \times 4 / 6 = 4\) 个八面体空隙。✔️

\begin{center}
\begin{minipage}[t]{0.48\textwidth}\centering
\includegraphics[width=0.95\linewidth,height=0.30\textheight,keepaspectratio]{puha-13-8b-tetrahedral-void.jpg}
\captionof{figure}{四面体空隙—4个等径球围成}
\end{minipage}
\hfill
\begin{minipage}[t]{0.48\textwidth}\centering
\includegraphics[width=0.95\linewidth,height=0.30\textheight,keepaspectratio]{puha-13-8c-octahedral-void.jpg}
\captionof{figure}{八面体空隙—6个等径球围成}
\end{minipage}
\end{center}

\begin{figure}[H]\centering
\includegraphics[width=0.55\textwidth,keepaspectratio]{puha-13-8a-close-packed-bilayer.jpg}
\caption{密置双层—第二层球占据第一层凹陷的一半位置}
\end{figure}

\begin{figure}[H]\centering
\includegraphics[width=0.55\textwidth,keepaspectratio]{puha-13-9a-ccp-stacking.jpg}
\caption{CCP(ABCABC…层序)}
\end{figure}

\begin{figure}[H]\centering
\includegraphics[width=0.55\textwidth,keepaspectratio]{puha-13-9a-hcp-stacking.jpg}
\caption{HCP(ABAB…层序)}
\end{figure}

\begin{quote}
相关知识点：密堆积 · 金属晶体 · 配位数 · 简单立方堆积 · 立方密堆积 · 六方密堆积 · 体心立方堆积
\end{quote}

\begin{center}\rule{0.5\linewidth}{0.5pt}\end{center}

\subsection{十五、离子晶体结构类型}\label{ux5341ux4e94ux79bbux5b50ux6676ux4f53ux7ed3ux6784ux7c7bux578b}

\begin{quote}
\textbf{关联知识点}：离子晶体 · 离子半径
\end{quote}

\subsubsection{15.1 为什么从金属堆积讲到离子晶体？}\label{ux4e3aux4ec0ux4e48ux4eceux91d1ux5c5eux5806ux79efux8bb2ux5230ux79bbux5b50ux6676ux4f53}

金属晶体是\textbf{等径球堆积}，但离子晶体不同---\textbf{负离子比正离子大得多}(通常2-3倍)。因此离子晶体的结构可以看作：

\begin{quote}
\textbf{核心图景}：负离子按一定方式堆积，正离子选择性填充在负离子围成的空隙中。
\end{quote}

这就是''堆积-填隙''模型，是理解离子晶体最直观的语言。

\begin{center}
\begin{minipage}[t]{0.48\textwidth}\centering
\includegraphics[width=0.95\linewidth,height=0.30\textheight,keepaspectratio]{puha-13-11a-nacl.jpg}
\captionof{figure}{NaCl型—Cl面心立方堆积，Na占全部八面体空隙(6:6)}
\end{minipage}
\hfill
\begin{minipage}[t]{0.48\textwidth}\centering
\includegraphics[width=0.95\linewidth,height=0.30\textheight,keepaspectratio]{puha-13-11b-cscl.jpg}
\captionof{figure}{CsCl型—Cl简单立方堆积，Cs占立方体空隙(8:8)}
\end{minipage}
\end{center}

\begin{figure}[H]\centering
\includegraphics[width=0.55\textwidth,keepaspectratio]{puha-13-11c-cubic-zns.jpg}
\caption{立方ZnS型—S面心立方堆积，Zn占半数四面体空隙(4:4)}
\end{figure}

{\def\LTcaptype{none} % do not increment counter
\begin{longtable}[]{@{}llc@{}}
\toprule\noalign{}
负离子堆积方式 & 围成的空隙类型 & 配位数 \\
\midrule\noalign{}
\endhead
\bottomrule\noalign{}
\endlastfoot
简单立方 & 立方体空隙 & 8 \\
面心立方(ccp/fcc) & 正四面体空隙、正八面体空隙 & 4 或 6 \\
六方(hcp) & 正四面体空隙、正八面体空隙 & 4 或 6 \\
\end{longtable}
}

\subsubsection{15.2 常见AB型离子晶体}\label{ux5e38ux89c1abux578bux79bbux5b50ux6676ux4f53}

\textbf{(1) NaCl型(岩盐型，6:6配位)}

\textbf{核心描述}：这是最常见、最基础的离子晶体结构。Cl\textsuperscript{-} 按面心立方最密堆积(ccp)，Na\textsuperscript{+} 占据全部正八面体空隙(100\%填隙率)。

{\def\LTcaptype{none} % do not increment counter
\begin{longtable}[]{@{}
  >{\raggedright\arraybackslash}p{(\linewidth - 2\tabcolsep) * \real{0.5000}}
  >{\raggedright\arraybackslash}p{(\linewidth - 2\tabcolsep) * \real{0.5000}}@{}}
\toprule\noalign{}
\begin{minipage}[b]{\linewidth}\raggedright
特征
\end{minipage} & \begin{minipage}[b]{\linewidth}\raggedright
描述
\end{minipage} \\
\midrule\noalign{}
\endhead
\bottomrule\noalign{}
\endlastfoot
负离子堆积 & Cl\textsuperscript{-} 按面心立方最密堆积(ccp/fcc) \\
正离子位置 & Na\textsuperscript{+} 占据\textbf{全部}正八面体空隙(100\%填隙率) \\
配位比 & 6:6(每个 Na\textsuperscript{+} 被 6 个 Cl\textsuperscript{-} 包围，每个 Cl\textsuperscript{-} 被 6 个 Na\textsuperscript{+} 包围) \\
晶胞内含 & 4 个 NaCl \\
原子坐标 & Cl\textsuperscript{-}(0,0,0; ½,½,0; ½,0,½; 0,½,½)；Na\textsuperscript{+}(½,½,½; 0,0,½; ½,0,0; 0,½,0) \\
实例 & NaF、KCl、MgO、CaO、NiO \\
\end{longtable}
}

\textbf{NaCl晶胞的原子计数}：
- Cl\textsuperscript{-}：顶点(\(8 \times 1/8 = 1\))+ 面心(\(6 \times 1/2 = 3\))= \textbf{4个}
- Na\textsuperscript{+}：体心(\(1\))+ 棱心(\(12 \times 1/4 = 3\))= \textbf{4个}

\begin{figure}[H]\centering
\includegraphics[width=0.55\textwidth,keepaspectratio]{nacl-structure-unit-cell.jpg}
\caption{NaCl型晶胞结构}
\end{figure}

\begin{quote}
\textbf{认知冲突}：「NaCl型结构 \不等于 面心立方格子！」---NaCl是两套面心立方格子(Na\textsuperscript{+}一套、Cl\textsuperscript{-}一套)互相穿插。面心立方格子要求八个顶点和六个面心是同种原子，不是这样交替的。
\textbf{区分口诀}：``金属看格子(A1/A2/A3)，离子看类型(NaCl/CsCl/ZnS)''
\end{quote}

\textbf{(2) CsCl型(8:8配位)}

{\def\LTcaptype{none} % do not increment counter
\begin{longtable}[]{@{}ll@{}}
\toprule\noalign{}
特征 & 描述 \\
\midrule\noalign{}
\endhead
\bottomrule\noalign{}
\endlastfoot
负离子堆积 & Cl\textsuperscript{-} 按简单立方堆积 \\
正离子位置 & Cs\textsuperscript{+} 占据立方体空隙(体心) \\
配位比 & 8:8 \\
晶胞内含 & 1 个 CsCl \\
原子坐标 & Cl\textsuperscript{-}(0,0,0)；Cs\textsuperscript{+}(½,½,½) \\
实例 & CsCl、CsBr、NH\textsubscript{4}Cl \\
\end{longtable}
}

\begin{figure}[H]\centering
\includegraphics[width=0.55\textwidth,keepaspectratio]{cscl-structure-unit-cell.jpg}
\caption{CsCl型结构}
\end{figure}

\begin{quote}
\textbf{认知冲突}：「CsCl是不是体心立方结构？」(发生率\textasciitilde50\%)---不是！CsCl是\textbf{简单立方}(两套简单立方格子互相穿插)。体心立方的顶点和体心是同种原子(如\(\alpha\)-Fe)，但CsCl的顶点是Cl\textsuperscript{-}、体心是Cs\textsuperscript{+}。\textbf{同种离子各自单看}，都是简单立方格子。
\end{quote}

\begin{quote}
\textbf{易错提醒}：CsCl型的点阵类型是\textbf{简单立方P}(不是体心立方I)！因为顶点和体心是不同种类的原子。
\end{quote}

\textbf{(3) 立方ZnS型(闪锌矿型，4:4配位)}

{\def\LTcaptype{none} % do not increment counter
\begin{longtable}[]{@{}
  >{\raggedright\arraybackslash}p{(\linewidth - 2\tabcolsep) * \real{0.5000}}
  >{\raggedright\arraybackslash}p{(\linewidth - 2\tabcolsep) * \real{0.5000}}@{}}
\toprule\noalign{}
\begin{minipage}[b]{\linewidth}\raggedright
特征
\end{minipage} & \begin{minipage}[b]{\linewidth}\raggedright
描述
\end{minipage} \\
\midrule\noalign{}
\endhead
\bottomrule\noalign{}
\endlastfoot
负离子堆积 & S\textsuperscript{2-} 按面心立方最密堆积 \\
正离子位置 & Zn\textsuperscript{2+} 占据\textbf{一半}正四面体空隙(50\%填隙率) \\
配位比 & 4:4(Zn\textsuperscript{2+} 位于正四面体中心，配位4个S\textsuperscript{2-}) \\
晶胞内含 & 4 个 ZnS \\
原子坐标 & S\textsuperscript{2-}(0,0,0; ½,½,0; ½,0,½; 0,½,½)；Zn\textsuperscript{2+}(¼,¼,¼; ¾,¾,¼; ¾,¼,¾; ¼,¾,¾) \\
实例 & BeS、CuCl、GaAs、CdS \\
\end{longtable}
}

\begin{figure}[H]\centering
\includegraphics[width=0.55\textwidth,keepaspectratio]{zns-structure-unit-cell.jpg}
\caption{立方ZnS型(闪锌矿)结构}
\end{figure}

\textbf{(4) 六方ZnS型(纤锌矿型，4:4配位)}

{\def\LTcaptype{none} % do not increment counter
\begin{longtable}[]{@{}ll@{}}
\toprule\noalign{}
特征 & 描述 \\
\midrule\noalign{}
\endhead
\bottomrule\noalign{}
\endlastfoot
负离子堆积 & S\textsuperscript{2-} 按六方最密堆积 \\
正离子位置 & Zn\textsuperscript{2+} 占据一半正四面体空隙 \\
配位比 & 4:4 \\
晶胞内含 & 2 个 ZnS \\
实例 & ZnO、BeO、GaN、AlN \\
\end{longtable}
}

\begin{figure}[H]\centering
\includegraphics[width=0.55\textwidth,keepaspectratio]{ZnS.png}
\caption{六方ZnS型结构(纤锌矿)}
\end{figure}

\subsubsection{15.3 其他常见结构}\label{ux5176ux4ed6ux5e38ux89c1ux7ed3ux6784}

\textbf{CaF\textsubscript{2}型(萤石型，8:4配位)}

\begin{itemize}
\tightlist
\item
  Ca\textsuperscript{2+} 按面心立方排列
\item
  F\textsuperscript{-} 占据 Ca\textsuperscript{2+} 围成的\textbf{全部}正四面体空隙(8个)
\item
  配位比：Ca\textsuperscript{2+} 为 8 配位(不规则立方体配位)，F\textsuperscript{-} 为 4 配位(正四面体)
\item
  实例：CaF\textsubscript{2}、BaF\textsubscript{2}、SrF\textsubscript{2}、UO\textsubscript{2}
\item
  \textbf{反萤石型}：正负离子位置互换---如 Li\textsubscript{2}O、Na\textsubscript{2}O(O\textsuperscript{2-} 按面心立方排列，Li\textsuperscript{+}/Na\textsuperscript{+} 填全部四面体空隙)
\end{itemize}

\begin{figure}[H]\centering
\includegraphics[width=0.55\textwidth,keepaspectratio]{CaF2.png}
\caption{CaF萤石型结构}
\end{figure}

\textbf{TiO\textsubscript{2}型(金红石型，6:3配位)}

\begin{itemize}
\tightlist
\item
  Ti\textsuperscript{4+} 位于晶胞顶点和体心(四方晶系)
\item
  O\textsuperscript{2-} 位于 TiO\textsubscript{6} 八面体之间的共顶点连接
\item
  配位比：Ti\textsuperscript{4+} 为 6 配位(八面体)，O\textsuperscript{2-} 为 3 配位(平面三角形配位)
\item
  实例：TiO\textsubscript{2}、MnO\textsubscript{2}、PbO\textsubscript{2}、SnO\textsubscript{2}
\end{itemize}

\begin{figure}[H]\centering
\includegraphics[width=0.55\textwidth,keepaspectratio]{TiO2.jpg}
\caption{金红石型TiO结构—Ti位于四方晶胞顶点和体心，O形成八面体配位}
\end{figure}

\textbf{(4) 钙钛矿型(ABO\textsubscript{3}，复杂配位)}

钙钛矿是另一类重要离子晶体结构，通式 ABO\textsubscript{3}。A位为大阳离子(如 Ca\textsuperscript{2+}、Ba\textsuperscript{2+})，B位为小阳离子(如 Ti\textsuperscript{4+}、Zr\textsuperscript{4+})，O\textsuperscript{2-} 位于面心位置形成八面体网络。

\begin{figure}[H]\centering
\includegraphics[width=0.55\textwidth,keepaspectratio]{perovskite-full.jpg}
\caption{钙钛矿(CaTiO)型结构—蓝色为B位Ti(八面体配位)，大灰色为A位Ca，红色为O}
\end{figure}

\subsubsection{15.4 结构类型对比总表}\label{ux7ed3ux6784ux7c7bux578bux5bf9ux6bd4ux603bux8868}

{\def\LTcaptype{none} % do not increment counter
\begin{longtable}[]{@{}
  >{\raggedright\arraybackslash}p{(\linewidth - 12\tabcolsep) * \real{0.1290}}
  >{\raggedright\arraybackslash}p{(\linewidth - 12\tabcolsep) * \real{0.1290}}
  >{\raggedright\arraybackslash}p{(\linewidth - 12\tabcolsep) * \real{0.1290}}
  >{\raggedright\arraybackslash}p{(\linewidth - 12\tabcolsep) * \real{0.1290}}
  >{\centering\arraybackslash}p{(\linewidth - 12\tabcolsep) * \real{0.1613}}
  >{\centering\arraybackslash}p{(\linewidth - 12\tabcolsep) * \real{0.1613}}
  >{\centering\arraybackslash}p{(\linewidth - 12\tabcolsep) * \real{0.1613}}@{}}
\toprule\noalign{}
\begin{minipage}[b]{\linewidth}\raggedright
结构类型
\end{minipage} & \begin{minipage}[b]{\linewidth}\raggedright
堆积离子
\end{minipage} & \begin{minipage}[b]{\linewidth}\raggedright
堆积方式
\end{minipage} & \begin{minipage}[b]{\linewidth}\raggedright
填隙离子
\end{minipage} & \begin{minipage}[b]{\linewidth}\centering
空隙类型
\end{minipage} & \begin{minipage}[b]{\linewidth}\centering
填隙率
\end{minipage} & \begin{minipage}[b]{\linewidth}\centering
配位比
\end{minipage} \\
\midrule\noalign{}
\endhead
\bottomrule\noalign{}
\endlastfoot
\textbf{CsCl} & Cl\textsuperscript{-} & 简单立方 & Cs\textsuperscript{+} & 立方体 & 100\% & 8:8 \\
\textbf{NaCl} & Cl\textsuperscript{-} & 面心立方(ccp) & Na\textsuperscript{+} & 正八面体 & 100\% & 6:6 \\
\textbf{立方ZnS} & S\textsuperscript{2-} & 面心立方(ccp) & Zn\textsuperscript{2+} & 正四面体 & 50\% & 4:4 \\
\textbf{六方ZnS} & S\textsuperscript{2-} & 六方(hcp) & Zn\textsuperscript{2+} & 正四面体 & 50\% & 4:4 \\
\textbf{CaF\textsubscript{2}} & Ca\textsuperscript{2+} & 面心立方(ccp) & F\textsuperscript{-} & 正四面体 & 100\% & 8:4 \\
\textbf{反萤石} & O\textsuperscript{2-} & 面心立方(ccp) & Li\textsuperscript{+} & 正四面体 & 100\% & 4:8 \\
\textbf{TiO\textsubscript{2}} & Ti\textsuperscript{4+} & 四方体心 & O\textsuperscript{2-} & 八面体链 & --- & 6:3 \\
\end{longtable}
}

\begin{quote}
\textbf{练一练}：填写下表：
\end{quote}

{\def\LTcaptype{none} % do not increment counter
\begin{longtable}[]{@{}lccc@{}}
\toprule\noalign{}
结构类型 & 配位比 & 正离子空隙类型 & 填隙率 \\
\midrule\noalign{}
\endhead
\bottomrule\noalign{}
\endlastfoot
NaCl & ( ) & ( ) & ( ) \\
CsCl & ( ) & ( ) & ( ) \\
立方ZnS & ( ) & ( ) & ( ) \\
CaF\textsubscript{2} & ( ) & ( ) & ( ) \\
\end{longtable}
}

(答案见§二十二参考答案)

\begin{quote}
相关 知识点：NaCl型结构 · CsCl型结构 · 立方ZnS型结构 · 六方ZnS型结构 · CaF2型结构 · NiAs型结构 · 反萤石型结构 · 金红石型结构
\end{quote}

\begin{center}\rule{0.5\linewidth}{0.5pt}\end{center}

\subsection{十六、半径比规则}\label{ux5341ux516dux534aux5f84ux6bd4ux89c4ux5219}

\begin{quote}
\textbf{关联知识点}：离子半径 · 半径比规则
\end{quote}

\subsubsection{16.1 离子晶体的''尺寸决定论''}\label{ux79bbux5b50ux6676ux4f53ux7684ux5c3aux5bf8ux51b3ux5b9aux8bba}

正离子填充什么类型的空隙---是四面体(4配位)、八面体(6配位)还是立方体(8配位)---这由什么决定？直觉告诉我们，\textbf{如果正离子太小，就填不进大的空隙；太大又会撑开负离子骨架}。半径比规则正是这一朴素直觉的量化表达。

\subsubsection{16.2 临界值几何推导}\label{ux4e34ux754cux503cux51e0ux4f55ux63a8ux5bfc}

半径比规则的核心思想：在纯离子模型中，正、负离子之间的静电吸引是最主要作用力。如果正离子太小，负离子外壳层就会互相接触---同种电荷相斥，结构不稳定。\textbf{最稳定的结构恰好是正离子刚好能接触到所有周围负离子的状态}---这就是''临界半径比''的几何原因。

\textbf{四配位(正四面体空隙)：\(r_+/r_- \geq 0.225\)}

正四面体的几何关系：四个球在顶点，中心恰好放入一个刚好接触四个球的球。

设负离子半径 \(r_-\)，正离子半径 \(r_+\)，正四面体棱边长 \(a = 2r_-\)，体心到顶点距离 \(d = \frac{\sqrt{6}}{4}a = \frac{\sqrt{6}}{4} \cdot 2r_- = \frac{\sqrt{6}}{2}r_-\)

临界条件：\(d = r_- + r_+\) \箭 \(r_+ = d - r_- = \left(\frac{\sqrt{6}}{2} - 1\right)r_-\) \箭 \(\boxed{\frac{r_+}{r_-} = \frac{\sqrt{6}}{2} - 1 \approx 0.225}\)

\textbf{六配位(正八面体空隙)：\(r_+/r_- \geq 0.414\)}

八面体xy平面截面为正方形：四个负离子在正方形顶点，正离子在中心。

正方形边长 \(a = 2r_-\)，对角线长 \(\sqrt{2}a = 2\sqrt{2}r_-\)，中心到顶点距离 \(= \frac{\sqrt{2}a}{2} = \sqrt{2}r_-\)

临界条件：\(\sqrt{2}r_- = r_- + r_+\) \箭 \(r_+ = (\sqrt{2} - 1)r_-\) \箭 \(\boxed{\frac{r_+}{r_-} = \sqrt{2} - 1 \approx 0.414}\)

\begin{figure}[H]\centering
\includegraphics[width=0.55\textwidth,keepaspectratio]{图-八面体空隙截面-半径比临界条件.png}
\caption{八面体空隙截面—4个负离子(灰)在正方形顶点，正离子(黑)在中心，临界条件 $r_+/r_- = \sqrt{2}-1 \approx 0.414$}
\end{figure}

\textbf{八配位(立方体空隙)：\(r_+/r_- \geq 0.732\)}

立方体体对角线 = \(\sqrt{3}a\)，棱长 \(a = 2r_-\)，体对角线长 \(= 2\sqrt{3}r_-\)

中心到顶点距离 = \(\frac{2\sqrt{3}r_-}{2} = \sqrt{3}r_-\)

临界条件：\(\sqrt{3}r_- = r_- + r_+\) \箭 \(r_+ = (\sqrt{3} - 1)r_-\) \箭 \(\boxed{\frac{r_+}{r_-} = \sqrt{3} - 1 \approx 0.732}\)

\begin{figure}[H]\centering
\includegraphics[width=0.55\textwidth,keepaspectratio]{radius_ratio_decision.png}
\end{figure}

\subsubsection{16.3 半径比与配位数的关系}\label{ux534aux5f84ux6bd4ux4e0eux914dux4f4dux6570ux7684ux5173ux7cfb}

{\def\LTcaptype{none} % do not increment counter
\begin{longtable}[]{@{}cccl@{}}
\toprule\noalign{}
\(r_+/r_-\) 范围 & 可能配位数 & 空隙类型 & 典型结构 \\
\midrule\noalign{}
\endhead
\bottomrule\noalign{}
\endlastfoot
0.155 \textasciitilde{} 0.225 & 3 & 三角形 & --- \\
0.225 \textasciitilde{} 0.414 & 4 & 正四面体 & ZnS 型 \\
0.414 \textasciitilde{} 0.732 & 6 & 正八面体 & NaCl 型 \\
0.732 \textasciitilde{} 1.0 & 8 & 立方体 & CsCl 型 \\
\大于等于 1.0 & 12 & 金属堆积 & --- \\
\end{longtable}
}

\begin{quote}
\textbf{认知冲突}：「半径比规则是几何必要条件，不是充分条件」---rc/ra 在范围内只说明几何上允许这个配位数，但实际结构还受离子极化、共价成分、温度压力等因素影响。\textbf{AgI 的 rc/ra \约等于 0.51(NaCl型范围)，但实际是ZnS型(4配位)}---因为Ag-I键有明显共价成分。
\end{quote}

\begin{quote}
\textbf{易错提醒}：半径比规则给出的是\textbf{最小}半径比。若正离子半径比该临界值更小，正、负离子不能接触，结构不稳定。但半径比规则并非绝对---当离子极化显著时(共价成分增多)，该规则的适用性下降。
\end{quote}

\begin{quote}
\textbf{练一练}：已知离子半径(单位pm)：K\textsuperscript{+}=138, Ca\textsuperscript{2+}=100, Na\textsuperscript{+}=102, O\textsuperscript{2-}=140, Cl\textsuperscript{-}=181。判断：(1) KCl 的预期结构类型；(2) CaO 的预期结构类型。
(答案见§二十二参考答案)
\end{quote}

\begin{quote}
相关知识点：离子半径 · 半径比规则
\end{quote}

\begin{center}\rule{0.5\linewidth}{0.5pt}\end{center}

\subsection{十七、Pauling 规则与离子极化}\label{ux5341ux4e03pauling-ux89c4ux5219ux4e0eux79bbux5b50ux6781ux5316}

\begin{quote}
\textbf{关联知识点}：Pauling规则 · 离子极化 · 晶格能
\end{quote}

\subsubsection{17.1 Pauling 电价规则---离子晶体结构的''电荷守恒''}\label{pauling-ux7535ux4ef7ux89c4ux5219ux79bbux5b50ux6676ux4f53ux7ed3ux6784ux7684ux7535ux8377ux5b88ux6052}

\textbf{问题}：在一个离子晶体中，正离子周围的负离子怎么排布才是最稳定的？Pauling 的第一规则(电价规则)给出了漂亮的答案：

\begin{quote}
\textbf{Pauling电价规则}：在稳定的离子晶体中，每个负离子的电价等于或近似等于从邻近正离子分配给该负离子的静电键强度总和。
\end{quote}

\[
Z_- = \sum_i \frac{Z_{+,i}}{\text{CN}_i}
\]

其中 \(Z_-\) 为负离子的电荷，\(Z_{+,i}\) 为第 i 个正离子的电荷，\(\text{CN}_i\) 为正离子的配位数。

\begin{quote}
\textbf{课堂原话}：``Pauling电价规则本质上是一本\textbf{能量账本}---每个正离子把自己'电力'平均分给所有邻居，所有邻居加起来必须刚好还完负离子的'债'。''
\end{quote}

\textbf{应用示例---判断NaCl型结构的合理性}：

NaCl型中，Cl\textsuperscript{-} 被 6 个 Na\textsuperscript{+} 包围(八面体配位)：
- 每个 Na\textsuperscript{+} 分配给 Cl\textsuperscript{-} 的静电键强度 = \(1/6\)
- 6 个 Na\textsuperscript{+} 总贡献 = \(6 \times 1/6 = 1\) = Cl\textsuperscript{-} 的电价 \箭 **结构合理** {[}完成{]}

\textbf{应用示例---判断CaF\textsubscript{2}型结构的合理性}：

Ca\textsuperscript{2+} 四面体配位(CN=8)，F\textsuperscript{-} 四面体配位(CN=4)：
- 每个 Ca\textsuperscript{2+} 分配给 F\textsuperscript{-} 的静电键强度 = \(2/8 = 1/4\)
- 4 个 Ca\textsuperscript{2+} 对 F\textsuperscript{-} 总贡献 = \(4 \times 1/4 = 1\) = F\textsuperscript{-} 的电价 \箭 **结构合理** {[}完成{]}

\begin{figure}[H]\centering
\includegraphics[width=0.55\textwidth,keepaspectratio]{图-八面体配位立体图.png}
\caption{八面体配位—6个负离子(灰)在八面体顶点，正离子(黑)在中心，配位数 = 6}
\end{figure}

\begin{quote}
\textbf{练一练}：采用 Pauling 电价规则检验 MgO(NaCl 型)的结构合理性。
(答案见§二十二参考答案)
\end{quote}

\subsubsection{17.2 Pauling 其他规则}\label{pauling-ux5176ux4ed6ux89c4ux5219}

\begin{itemize}
\tightlist
\item
  \textbf{第二规则(共顶点/共棱/共面规则)}：配位多面体倾向于共顶点连接，共棱和共面连接会降低稳定性(正离子间距减小，排斥增大)

  \begin{itemize}
  \tightlist
  \item
    共顶点 \箭 正离子间距最大，排斥最小
  \item
    共棱 \箭 中等
  \item
    共面 \箭 正离子间距最小，排斥最大
  \end{itemize}
\item
  \textbf{第三规则(规则性)}：配位多面体倾向于尽可能规则
\end{itemize}

\subsubsection{17.3 离子极化与 Fajans 规则}\label{ux79bbux5b50ux6781ux5316ux4e0e-fajans-ux89c4ux5219}

\textbf{为什么有些离子晶体''不像离子晶体''？}

离子极化：\textbf{正离子使负离子电子云变形(极化作用)，负离子被极化(变形性)}。当极化足够强时，键从离子型向共价型过渡---配位数下降、物性改变。

\begin{figure}[H]\centering
\includegraphics[width=0.55\textwidth,keepaspectratio]{puha-13-17a-no-polarization.jpg}
\caption{(a) 理想离子键(无极化)—电子云球形对称}
\end{figure}

\begin{figure}[H]\centering
\includegraphics[width=0.55\textwidth,keepaspectratio]{puha-13-17b-slight-polarization.jpg}
\caption{(b) 轻微极化—负离子电子云略变形}
\end{figure}

\begin{figure}[H]\centering
\includegraphics[width=0.55\textwidth,keepaspectratio]{puha-13-17c-strong-polarization.jpg}
\caption{(c) 较强极化—电子云明显偏向正离子}
\end{figure}

\begin{figure}[H]\centering
\includegraphics[width=0.55\textwidth,keepaspectratio]{puha-13-17d-covalent-bond.jpg}
\caption{(d) 强烈极化—电子云严重重叠}
\end{figure}

{\def\LTcaptype{none} % do not increment counter
\begin{longtable}[]{@{}lll@{}}
\toprule\noalign{}
因素 & 极化作用(正离子) & 变形性(负离子) \\
\midrule\noalign{}
\endhead
\bottomrule\noalign{}
\endlastfoot
电荷 & 电荷越高 \箭 极化越强 & 电荷越高 \箭 变形性越大 \\
半径 & 半径越小 \箭 极化越强 & 半径越大 \箭 变形性越大 \\
电子构型 & 18e\textsuperscript{-}/18+2e\textsuperscript{-} \textgreater{} 9\textasciitilde17e\textsuperscript{-} \textgreater{} 8e\textsuperscript{-} & 阴离子通常 \textgreater{} 阳离子 \\
\end{longtable}
}

\begin{quote}
\textbf{认知冲突}：「为什么AgI的半径比(0.523)落在NaCl型范围，实际却是ZnS型？」---Ag\textsuperscript{+}是18电子构型(4d\textsuperscript{10})，极化作用强；I\textsuperscript{-}半径大(220pm)，变形性大。Ag\textsuperscript{+}的强极化使电子云严重畸变\箭键的共价性增大\箭配位数从6降为4。
\end{quote}

\textbf{经典反例}：

{\def\LTcaptype{none} % do not increment counter
\begin{longtable}[]{@{}
  >{\raggedright\arraybackslash}p{(\linewidth - 6\tabcolsep) * \real{0.2222}}
  >{\centering\arraybackslash}p{(\linewidth - 6\tabcolsep) * \real{0.2778}}
  >{\centering\arraybackslash}p{(\linewidth - 6\tabcolsep) * \real{0.2778}}
  >{\raggedright\arraybackslash}p{(\linewidth - 6\tabcolsep) * \real{0.2222}}@{}}
\toprule\noalign{}
\begin{minipage}[b]{\linewidth}\raggedright
化合物
\end{minipage} & \begin{minipage}[b]{\linewidth}\centering
预测结构(半径比规则)
\end{minipage} & \begin{minipage}[b]{\linewidth}\centering
实际结构
\end{minipage} & \begin{minipage}[b]{\linewidth}\raggedright
解释
\end{minipage} \\
\midrule\noalign{}
\endhead
\bottomrule\noalign{}
\endlastfoot
AgI & \(r_+/r_- = 115/220 = 0.523\) \箭 NaCl 型(6配位) & ZnS 型(4配位) & Ag\textsuperscript{+} 为 18e\textsuperscript{-} 构型，极化作用强；I\textsuperscript{-} 半径大 \箭 共价性增强 \箭 配位数降低 \\
CuCl & \(r_+/r_- > 0.414\) \箭 NaCl 型 & ZnS 型 & Cu\textsuperscript{+} 为 18e\textsuperscript{-} 构型，极化导致结构转变 \\
Al\textsubscript{2}O\textsubscript{3}(刚玉) & 预期较高配位 & 实际CN=6(八面体配位) & Al\textsuperscript{3+} 高电荷 + 小半径 \箭 极化力极强 \\
\end{longtable}
}

\begin{quote}
\textbf{练一练}：解释为什么 BeO 采取六方 ZnS 型(4配位)而不是 NaCl 型(6配位)。已知 Be\textsuperscript{2+}半径31pm，O\textsuperscript{2-}半径140pm。
(答案见§二十二参考答案)
\end{quote}

\begin{quote}
相关 KP：Pauling规则 · 离子极化 · 晶格能 · 离子键
\end{quote}

\begin{center}\rule{0.5\linewidth}{0.5pt}\end{center}

\subsection{十八、其他晶体类型}\label{ux5341ux516bux5176ux4ed6ux6676ux4f53ux7c7bux578b}

\begin{quote}
\textbf{关联知识点}：原子晶体 · 分子晶体 · 金属晶体 · 金属键
\end{quote}

\subsubsection{18.1 四大晶体类型对比}\label{ux56dbux5927ux6676ux4f53ux7c7bux578bux5bf9ux6bd4}

{\def\LTcaptype{none} % do not increment counter
\begin{longtable}[]{@{}
  >{\raggedright\arraybackslash}p{(\linewidth - 8\tabcolsep) * \real{0.2000}}
  >{\raggedright\arraybackslash}p{(\linewidth - 8\tabcolsep) * \real{0.2000}}
  >{\raggedright\arraybackslash}p{(\linewidth - 8\tabcolsep) * \real{0.2000}}
  >{\raggedright\arraybackslash}p{(\linewidth - 8\tabcolsep) * \real{0.2000}}
  >{\raggedright\arraybackslash}p{(\linewidth - 8\tabcolsep) * \real{0.2000}}@{}}
\toprule\noalign{}
\begin{minipage}[b]{\linewidth}\raggedright
晶体类型
\end{minipage} & \begin{minipage}[b]{\linewidth}\raggedright
结构单元
\end{minipage} & \begin{minipage}[b]{\linewidth}\raggedright
结合力
\end{minipage} & \begin{minipage}[b]{\linewidth}\raggedright
特点
\end{minipage} & \begin{minipage}[b]{\linewidth}\raggedright
实例
\end{minipage} \\
\midrule\noalign{}
\endhead
\bottomrule\noalign{}
\endlastfoot
\textbf{离子晶体} & 正负离子 & 离子键(静电作用) & 硬而脆、熔沸点高、水溶液导电、熔融导电 & NaCl、MgO、CaF\textsubscript{2} \\
\textbf{原子晶体(共价晶体)} & 原子 & 共价键 & 最硬、熔点最高、不导电(除少数)、耐腐蚀 & 金刚石、SiC、SiO\textsubscript{2} \\
\textbf{分子晶体} & 分子/原子 & 分子间力/氢键 & 软、熔沸点低、固态不导电 & 干冰、I\textsubscript{2}、冰、C\textsubscript{60} \\
\textbf{金属晶体} & 原子 & 金属键 & 导电导热、延展性好、有光泽 & Cu、Fe、Na、Al \\
\end{longtable}
}

\begin{figure}[H]\centering
\includegraphics[width=0.55\textwidth,keepaspectratio]{13-14-co2-molecular-crystal-unit-cell.jpg}
\caption{干冰(CO)分子晶体晶胞}
\end{figure}

\subsubsection{18.2 混合型晶体 --- 石墨}\label{ux6df7ux5408ux578bux6676ux4f53-ux77f3ux58a8}

石墨为\textbf{层状结构}：
- 层内：C 原子 sp\textsuperscript{2} 杂化，形成离域\派键(\派ₙᵐ)，键能大 \箭 **高熔点**
- 层间：范德华力结合 \箭 **易滑动**(润滑性)
- 离域\派电子可在层内自由移动 \箭 **导电性**(各向异性：平行于层面导电性远好于垂直方向)

\textbf{各向异性的本质原因}：层内共价键(强方向性)+ 层间分子间力(弱非方向性)= 性质在不同方向上大不同。

\begin{figure}[H]\centering
\includegraphics[width=0.55\textwidth,keepaspectratio]{graphite-structure.jpg}
\caption{石墨的层状结构}
\end{figure}

\begin{figure}[H]\centering
\includegraphics[width=0.48\textwidth,height=0.32\textheight,keepaspectratio]{diamond-structure.jpg}
\hspace{0.015\textwidth}
\includegraphics[width=0.48\textwidth,height=0.32\textheight,keepaspectratio]{graphite-structure.jpg}
\caption{金刚石(左)与石墨(右)的晶体结构对比—金刚石为三维共价网络(sp)，石墨为层状结构(层内sp共价键+层间范德华力)}
\end{figure}

\begin{quote}
相关 KP：原子晶体 · 分子晶体 · 金属晶体 · 金属键 · 离子键 · 石墨结构 · 金刚石型结构 · 碳的同素异形体 · 二氧化硅结构 · 原子坐标参数 · 分数坐标
\end{quote}

\begin{center}\rule{0.5\linewidth}{0.5pt}\end{center}

\subsection{十九、典型例题}\label{ux5341ux4e5dux5178ux578bux4f8bux9898}

\textbf{题目}：金属铜(Cu)采用面心立方堆积，已知 Cu 的原子半径为 128 pm，求晶胞参数 a 和晶体密度(Cu 的原子量 = 63.55 g/mol)。

\textbf{思路分析}：
1. 面心立方中面对角线方向原子接触 \箭 \(\sqrt{2}a = 4r\)
2. 晶胞含 4 个 Cu 原子
3. 密度 \(\rho = n \times M / (N_A \times a^3)\)

(答案见§二十二)

\begin{center}\rule{0.5\linewidth}{0.5pt}\end{center}

\textbf{题目}：已知离子半径如下(单位pm)：

{\def\LTcaptype{none} % do not increment counter
\begin{longtable}[]{@{}
  >{\centering\arraybackslash}p{(\linewidth - 18\tabcolsep) * \real{0.1000}}
  >{\centering\arraybackslash}p{(\linewidth - 18\tabcolsep) * \real{0.1000}}
  >{\centering\arraybackslash}p{(\linewidth - 18\tabcolsep) * \real{0.1000}}
  >{\centering\arraybackslash}p{(\linewidth - 18\tabcolsep) * \real{0.1000}}
  >{\centering\arraybackslash}p{(\linewidth - 18\tabcolsep) * \real{0.1000}}
  >{\centering\arraybackslash}p{(\linewidth - 18\tabcolsep) * \real{0.1000}}
  >{\centering\arraybackslash}p{(\linewidth - 18\tabcolsep) * \real{0.1000}}
  >{\centering\arraybackslash}p{(\linewidth - 18\tabcolsep) * \real{0.1000}}
  >{\centering\arraybackslash}p{(\linewidth - 18\tabcolsep) * \real{0.1000}}
  >{\centering\arraybackslash}p{(\linewidth - 18\tabcolsep) * \real{0.1000}}@{}}
\toprule\noalign{}
\begin{minipage}[b]{\linewidth}\centering
离子
\end{minipage} & \begin{minipage}[b]{\linewidth}\centering
K\textsuperscript{+}
\end{minipage} & \begin{minipage}[b]{\linewidth}\centering
Cs\textsuperscript{+}
\end{minipage} & \begin{minipage}[b]{\linewidth}\centering
Ca\textsuperscript{2+}
\end{minipage} & \begin{minipage}[b]{\linewidth}\centering
Na\textsuperscript{+}
\end{minipage} & \begin{minipage}[b]{\linewidth}\centering
Cl\textsuperscript{-}
\end{minipage} & \begin{minipage}[b]{\linewidth}\centering
Br\textsuperscript{-}
\end{minipage} & \begin{minipage}[b]{\linewidth}\centering
I\textsuperscript{-}
\end{minipage} & \begin{minipage}[b]{\linewidth}\centering
O\textsuperscript{2-}
\end{minipage} & \begin{minipage}[b]{\linewidth}\centering
S\textsuperscript{2-}
\end{minipage} \\
\midrule\noalign{}
\endhead
\bottomrule\noalign{}
\endlastfoot
半径 & 138 & 167 & 100 & 102 & 181 & 196 & 220 & 140 & 184 \\
\end{longtable}
}

判断下列化合物的晶体结构类型：(1) KBr；(2) CaO；(3) CsCl。

\textbf{思路分析}：
1. 计算 \(r_+/r_-\)
2. 对照半径比规则判断配位数和结构类型

(答案见§二十二)

\begin{center}\rule{0.5\linewidth}{0.5pt}\end{center}

\textbf{题目}：已知AgI中离子半径 Ag\textsuperscript{+}=115pm, I\textsuperscript{-}=220pm，按半径比规则应预测为哪种结构？实际AgI是哪种结构？为什么？

\textbf{思路分析}：
1. 先按半径比规则预测理论结构
2. 再考虑Ag\textsuperscript{+}的电子构型和I\textsuperscript{-}的变形性
3. 用离子极化理论解释偏差

(答案见§二十二)

\begin{center}\rule{0.5\linewidth}{0.5pt}\end{center}

\textbf{题目}：说明金刚石、NaCl、干冰(CO\textsubscript{2})、铜四种晶体的键型、硬度、熔点差异，并解释原因。

(答案见§二十二)

\textbf{题目}：已知金属Al采用面心立方堆积，晶胞参数 a = 405 pm。计算：
(1) Al的原子半径
(2) 每个晶胞中的四面体空隙数
(3) 假设把所有的四面体空隙填满H原子(H原子半径为 37 pm)，H原子之间是否接触？

(答案见§二十二)

\begin{center}\rule{0.5\linewidth}{0.5pt}\end{center}

\textbf{题目}：CaF\textsubscript{2}(萤石)中 Ca\textsuperscript{2+} 半径为 100 pm，F\textsuperscript{-} 半径为 133 pm。
(1) 计算半径比，预测 Ca\textsuperscript{2+} 的配位数。
(2) 实际 CaF\textsubscript{2} 中 Ca\textsuperscript{2+} 的配位数和 F\textsuperscript{-} 的配位数各是多少？
(3) 一个 CaF\textsubscript{2} 晶胞含几个 Ca\textsuperscript{2+}？几个 F\textsuperscript{-}？

(答案见§二十二)

\begin{center}\rule{0.5\linewidth}{0.5pt}\end{center}

\subsection{二十、方法总结}\label{ux4e8cux5341ux65b9ux6cd5ux603bux7ed3}

\subsubsection{20.1 晶体结构分析流程图}\label{ux6676ux4f53ux7ed3ux6784ux5206ux6790ux6d41ux7a0bux56fe}

\begin{Shaded}
\begin{Highlighting}[]
\NormalTok{给定化学式}
\NormalTok{   ↓}
\NormalTok{Step 1: 判断晶体类型（离子/原子/分子/金属/混合）}
\NormalTok{   由组成元素和物性推断}
\NormalTok{   ↓}
\NormalTok{Step 2: 确定堆积方式（金属/离子晶体）}
\NormalTok{   AB 型 → 算半径比 → 对照范围 → 确定结构类型}
\NormalTok{   ↓}
\NormalTok{Step 3: 分析配位数和配位多面体}
\NormalTok{   NaCl: 6 配位八面体; CsCl: 8 配位立方体; ZnS: 4 配位四面体}
\NormalTok{   ↓}
\NormalTok{Step 4: 晶体密度计算}
\NormalTok{   ρ = (n × M) / (N\_A × a³)}
\NormalTok{   关键：n → 晶胞中化学式单位数；a → 晶胞参数}
\NormalTok{   ↓}
\NormalTok{Step 5: 稳定性判断（必要时）}
\NormalTok{   半径比规则 → Pauling 电价规则 → 离子极化修正}
\end{Highlighting}
\end{Shaded}

\subsubsection{20.2 核心速查表}\label{ux6838ux5fc3ux901fux67e5ux8868}

{\def\LTcaptype{none} % do not increment counter
\begin{longtable}[]{@{}
  >{\raggedright\arraybackslash}p{(\linewidth - 4\tabcolsep) * \real{0.3333}}
  >{\raggedright\arraybackslash}p{(\linewidth - 4\tabcolsep) * \real{0.3333}}
  >{\raggedright\arraybackslash}p{(\linewidth - 4\tabcolsep) * \real{0.3333}}@{}}
\toprule\noalign{}
\begin{minipage}[b]{\linewidth}\raggedright
概念
\end{minipage} & \begin{minipage}[b]{\linewidth}\raggedright
关键公式/规律
\end{minipage} & \begin{minipage}[b]{\linewidth}\raggedright
注意点
\end{minipage} \\
\midrule\noalign{}
\endhead
\bottomrule\noalign{}
\endlastfoot
\textbf{7大晶系} & 轴长+轴角关系 & 六方注意120°；晶系由对称性定，不是边长 \\
\textbf{ccp空间利用率} & \(\pi/3\sqrt{2} \approx 74.05\%\) & 面对角线方向原子接触 \\
\textbf{bcc空间利用率} & \(\sqrt{3}\pi/8 \approx 68.02\%\) & 体对角线方向原子接触 \\
\textbf{hcp空间利用率} & \(\pi/3\sqrt{2} \approx 74.05\%\)(与ccp相同) & 理想轴比 c/a = 1.633 \\
\textbf{半径比规则} & 0.225\箭4, 0.414\箭6, 0.732\箭8 & 极化显著时规则失效 \\
\textbf{Pauling电价规则} & \(Z_- = \sum Z_+/ \text{CN}\) & 配位数不能加错 \\
\textbf{离子极化顺序} & 18e\textsuperscript{-} 构型 \textgreater{} 9\textasciitilde17e\textsuperscript{-} \textgreater{} 8e\textsuperscript{-} & AgI是经典反例 \\
\textbf{晶体密度} & \(\rho = nM/N_A a^3\) & 晶胞原子数 n 别数错 \\
\end{longtable}
}

\subsubsection{20.3 结构类型判断流程图}\label{ux7ed3ux6784ux7c7bux578bux5224ux65adux6d41ux7a0bux56fe}

\begin{Shaded}
\begin{Highlighting}[]
\NormalTok{给定 AB 型离子晶体}
\NormalTok{   ↓}
\NormalTok{计算 r₊/r₋}
\NormalTok{   ↓}
\NormalTok{┌─ \textless{}0.225 → 配位数3或更低（少见）}
\NormalTok{├─ 0.225\textasciitilde{}0.414 → 4配位 → ZnS型}
\NormalTok{├─ 0.414\textasciitilde{}0.732 → 6配位 → NaCl型}
\NormalTok{├─ 0.732\textasciitilde{}1.0 → 8配位 → CsCl型}
\NormalTok{└─ \textgreater{}1.0 → 配位数12 → 金属堆积}
\NormalTok{   ↓}
\NormalTok{是否有强极化？(Ag⁺, Cu⁺, Zn²⁺, Cd²⁺等)}
\NormalTok{  → 是：配位数可能下降（→ ZnS型或更低）}
\NormalTok{  → 否：按半径比规则预测}
\end{Highlighting}
\end{Shaded}

\begin{center}\rule{0.5\linewidth}{0.5pt}\end{center}

\subsection{二十一、练习题}\label{ux4e8cux5341ux4e00ux7ec3ux4e60ux9898}

\begin{enumerate}
\def\labelenumi{\arabic{enumi}.}
\item
  说出立方晶系、四方晶系、六方晶系的轴长/轴角特征。
\item
  计算面心立方(fcc)一个晶胞中的原子数，已知原子位于顶点和面心。
\item
  已知 CsCl 型晶体中 Cl\textsuperscript{-} 和 Cs\textsuperscript{+} 的坐标分别是 (0,0,0) 和 (½,½,½)，计算一个晶胞含多少个 CsCl 化学式单位。
\item
  已知 Na\textsuperscript{+} 半径为 102 pm，Cl\textsuperscript{-} 半径为 181 pm，判断 NaCl 的晶体结构类型和配位数。
\item
  背诵半径比规则的三档临界值。
\item
  已知 Ca\textsuperscript{2+} 半径为 100 pm，F\textsuperscript{-} 半径为 133 pm，判断 CaF\textsubscript{2} 中 Ca\textsuperscript{2+} 的配位数。
\item
  解释为什么 AgI 采取 ZnS 型(4配位)而非 NaCl 型(6配位)。
\item
  已知 MgO(NaCl型)和 NaF(NaCl型)具有相同结构，但 MgO 的熔点(2852°C)远高于 NaF(993°C)，解释原因。
\item
  采用 Pauling 电价规则检验 MgO(NaCl 型)的结构合理性。
\item
  推导 hcp 的理想轴比 c/a。
\item
  fcc 晶胞中有多少个八面体空隙？多少个四面体空隙？验证''每球对应1个八面体空隙+2个四面体空隙''。
\item
  已知 Be\textsuperscript{2+}=31pm, O\textsuperscript{2-}=140pm，判断 BeO 的预期结构。如果实际 BeO 为 ZnS 型，为什么？
\item
  Cu 的晶体结构为面心立方，原子半径为 128 pm，原子量为 63.55，计算 Cu 的理论密度。若实验测得密度为 8.96 g/cm\textsuperscript{3}，计算 Cu 晶体中实际晶胞参数是多少？
\item
  NiAs 型结构(NiAs，红镍矿)是另一种常见 AB 型结构---As 按六方密堆积(hcp)，Ni 填满全部八面体空隙。写出 Ni 的配位数和 As 的配位数，并画出结构示意。
\item
  某 AB\textsubscript{2} 型离子晶体，A\textsuperscript{2+} 半径 110 pm，B\textsuperscript{-} 半径 200 pm。预测其可能的结构类型，画出晶胞示意并计算每个晶胞中的化学式单位数。
\end{enumerate}

\begin{center}\rule{0.5\linewidth}{0.5pt}\end{center}

\subsection{二十二、参考答案与提示}\label{ux4e8cux5341ux4e8cux53c2ux8003ux7b54ux6848ux4e0eux63d0ux793a}

\textbf{19.1} \(a = 362\ \mathrm{pm}\)，\(\rho = 8.90\ \mathrm{g/cm}^3\)。

\textbf{19.2} (1) KBr: 0.704\箭NaCl型；(2) CaO: 0.714\箭NaCl型；(3) CsCl: 0.923\箭CsCl型。

\textbf{19.3} 半径比0.523预测NaCl型，实际ZnS型(Ag\textsuperscript{+}极化修正)。

\textbf{19.4} 硬度：金刚石\textgreater NaCl\textgreater Cu\textgreater 干冰。原因：共价键\textgreater 离子键\textgreater 金属键\textgreater 分子间力。

\textbf{19.5} \(r=143\) pm；8个四面体空隙；H半径37pm\textgreater 临界32.2pm，几何可填入。

\textbf{19.6} \(r_+/r_-=0.752\)\箭配位数8；实际Ca\textsuperscript{2+}8配位，F\textsuperscript{-}4配位；晶胞含4个CaF\textsubscript{2}(Ca\textsuperscript{2+}4个+F\textsuperscript{-}8个)。

\begin{enumerate}
\def\labelenumi{\arabic{enumi}.}
\item
  六方晶系(\(\gamma=120°\), \(a=b\neq c\))。
\item
  hcp 理想 \(c/a = 2\sqrt{2/3} \approx 1.633\)。
\item
  NaCl: 6:6, 正八面体空隙, 100\%；CsCl: 8:8, 立方体空隙, 100\%；立方ZnS: 4:4, 正四面体空隙, 50\%；CaF\textsubscript{2}: 8:4, 正四面体空隙, 100\%。
\item
  \begin{enumerate}
  \def\labelenumii{(\arabic{enumii})}
  \tightlist
  \item
    KCl: \(r_+/r_-=0.762\) \箭 预测CsCl型(实际NaCl型，极化修正)。(2) CaO: \(r_+/r_-=0.714\) \箭 NaCl型。
  \end{enumerate}
\item
  MgO: \(6\times(2/6)=2=O^{2-}\)电价 \箭 结构合理。BeO: Be\textsuperscript{2+}极化力极强 \箭 实际ZnS型(4配位)。
\item
  立方: \(a=b=c,\alpha=\beta=\gamma=90°\)；四方: \(a=b\neq c,\alpha=\beta=\gamma=90°\)；六方: \(a=b\neq c,\alpha=\beta=90°,\gamma=120°\)。
\item
  顶点\(8\times1/8\)+面心\(6\times1/2\)=\textbf{4个}。
\item
  顶点8个Cl\textsuperscript{-}×1/8 + 体心1个Cs\textsuperscript{+} = \textbf{1个CsCl}。
\item
  \(r_+/r_-=0.564\) \箭 **NaCl型\textbf{，配位数}6**。
\item
  0.225(四面体)、0.414(八面体)、0.732(立方体)。
\item
  \(r_+/r_-=0.752>0.732\) \箭 Ca\textsuperscript{2+}配位数\textbf{8}(CaF\textsubscript{2}型，配位比8:4)。
\item
  Ag\textsuperscript{+}为18e\textsuperscript{-}构型，极化作用强；I\textsuperscript{-}半径大、变形性大 \箭 键共价性增大 \箭 配位数6\箭4(ZnS型)。
\item
  晶格能\(U\propto Z_+Z_-/r\)，MgO电荷\(\pm2\)远大于NaF的\(\pm1\) \箭 MgO晶格能远大 \箭 熔点更高。
\item
  MgO: \(6\times(2/6)=2=O^{2-}\)电价 \箭 **结构合理** {[}完成{]}。
\item
  \(c/a=2\sqrt{2/3}\approx1.633\)。正四面体高\(h=\sqrt{6}/3\times a\)，hcp含2层 \箭 \(c=4h\)。
\item
  fcc: \textbf{4个八面体空隙} + \textbf{8个四面体空隙}。验证：\(6\times4/6=4\)；\(8\times4/4=8\)。
\item
  Be\textsuperscript{2+}电荷高、半径极小 \箭 极化力极强 \箭 共价性显著 \箭 4配位被锁定(ZnS型)。
\item
  \(\rho_{th}=4\times63.55/(N_A\times a^3)\)，\(a=4r/\sqrt{2}=362\) pm \箭 \(\rho_{th}\approx8.90\) g/cm\textsuperscript{3}。由\(\rho_{exp}=8.96\)反推\(a_{exp}\)。
\item
  NiAs型: As按hcp堆积，Ni填全部八面体空隙。Ni配位数=6(八面体)，As配位数=6(三棱柱)。晶胞含1个NiAs。
\item
  \(r_+/r_-=0.55\)，0.414\textasciitilde0.732 \箭 配位数6。AB\textsubscript{2}型 \箭 **金红石型**(TiO\textsubscript{2}，配位比6:3)。CaF\textsubscript{2}型要求\(r_+/r_->0.732\)，不满足。
\item
  \(a=4r/\sqrt{2}=362\) pm，\(\rho=8.90\) g/cm\textsuperscript{3}。
\item
  KBr: 0.704\箭NaCl型；CaO: 0.714\箭NaCl型；CsCl: 0.923\箭CsCl型。
\item
  半径比0.523预测NaCl型，实际ZnS型(Ag\textsuperscript{+}极化修正)。
\item
  硬度：金刚石\textgreater NaCl\textgreater Cu\textgreater 干冰。原因：共价键\textgreater 离子键\textgreater 金属键\textgreater 分子间力。
\item
  \(r=143\) pm；8个四面体空隙；H半径37pm\textgreater 临界32.2pm，几何可填入。
\item
  \(r_+/r_-=0.752\)\箭配位数8；实际Ca\textsuperscript{2+}8配位，F\textsuperscript{-}4配位；晶胞含4个CaF\textsubscript{2}(Ca\textsuperscript{2+}4个+F\textsuperscript{-}8个)。
\end{enumerate}

\begin{center}\rule{0.5\linewidth}{0.5pt}\end{center}

\subsection{二十三、延伸阅读}\label{ux4e8cux5341ux4e09ux5ef6ux4f38ux9605ux8bfb}

\begin{itemize}
\tightlist
\item
  《普通化学原理(第4版)》第13章：晶体结构与固体化学基础
\item
  宋天佑《无机化学》§10-11 章
\item
  Weller 等《无机化学》相关晶体结构章节
\item
  《化学竞赛初赛讲义》第 7 讲：晶体结构
\end{itemize}

\begin{center}\rule{0.5\linewidth}{0.5pt}\end{center}

\end{document}
